% Modernised reading — fonds Alexandre Grothendieck, shelfmark 140-1, pages 1-36.
%
% This is an INTERPRETATION, not a transcription. It re-reads the pages in
% today's notation and vocabulary, states the hypotheses the manuscript leaves
% implicit, and drops the critical apparatus so that the mathematics reads
% straight through. Everything the brackets and underlines carried — a doubtful
% reading, a notation he used differently, a missing passage — moves into a
% footnote. It is held to being mathematically correct as it stands; where the
% manuscript is loose or mistaken the true statement appears and the footnote
% says what the page has. In any dispute of substance the transcription
% governs: whoever cites Grothendieck cites batch-01.fr.tex and batch-02.fr.tex.
%
% Scope: the folder was read WHOLE before any of this was written — both
% batches, pages 1-20 and 21-36 — and it is written as ONE file for the whole
% shelfmark. Pages with no transcription: 2, 8, 21, 23, 25, 27, 36 (typed or
% blank leaves, per the transcriptions); 4, 6, 10, 12 were skipped by the first
% sitting of batch 1 and not re-examined, so this reading cannot say whether
% they carry anything of his.
%
% The spine: drafts for the Longue Marche (1981) on ONE object — the
% Teichmüller group T_{1,1} (= SL(2,Z)) and its central extension ST_{1,1}
% (= the preimage of SL(2,Z) in the universal cover of SL(2,R), i.e. the braid
% group B_3) — approached through the orbifold (U_{0,3}, S_3) and its group
% Pi^D_{0,3} = <rho, sigma | rho^3 = sigma^2 = 1> (= PSL(2,Z)). Stations:
% central extensions of Pi^D_n = Z/n * Z/2 (3-5, 11-13, 18); the n = 3 case
% and its finite quotients (7, 9, 14); T_{1,1}, ST_{1,1} (15, 17, 19); the
% equivariant Picard group of U_{0,3} (22, 24, 29 top); stacks over M_{1,1}
% (26); the diagram of extensions (28, 30); extension classes and the
% character of the mu_6-cover (32, 34); the answer ST_{1,1} = SPi^D_{0,3}
% (33, 35). Two side threads: surfaces with boundary (20) and combinatorial
% maps (16, 29 bottom, 31).
%
% Notation fixed for the whole folder. Pi^D_n = <rho, sigma | rho^n = sigma^2
% = 1>; D_n its dihedral quotient (sigma rho)^2 = 1; Pi_n the kernel, free on
% rho_0..rho_{n-1} with rho_{n-1}...rho_0 = 1. The page gives the SAME name
% \widetilde\Pi^D_n to two different central extensions: page 3's
% <rho~, sigma~ | rho~^n = sigma~^2> is kept as \widetilde\Pi^D_n; page 11's
% fibre product Pi^D_n x_{Z/2n} Z is written S\Pi^D_n, the name page 33 gives
% it for n = 3. Page 9's Pi'_{0,3} (index-3 subgroup of Pi_{0,3}) keeps its
% name; page 28's « Pi'_{0,3} = Pi_{0,3} x {1, omega} » is written Gamma(2).
% A prime on SL(2,Z)', GL(2,Z)', M'_{11} is read as « modulo ±1 ».
%
% Substantive departures from the pages, each footnoted where it occurs:
%  - p. 1: the boxed formula holds only when alpha^{-1}u commutes with rho
%    (and v != 0); the middle term of the page's det bracket has the wrong
%    sign; det(alpha^{-1}u) = mu/delta holds directly.
%  - p. 5: the question « invariant by sigma~ ? » is left open by the page;
%    for n odd the reading answers yes (abelianisation mod 4), marked as ours.
%  - p. 7: the NB « lifts of rho have order 6 » is false (rho.omega has order
%    3); the extension D'_3 -> D_3 is non-split because the lifts of sigma
%    have order 4. GL(2,Z)/Pi_{0,3} is a semidirect product, not the page's
%    direct product D'_3 x {1, tau}.
%  - p. 9: the page's presentation of D'_3 drops sigma(rho) = rho^{-1}; E
%    (order 36) cannot be a quotient of SL(2,Z/6Z) (abelianisations Z/12 and
%    Z/6); E is identified with page 28's D'''_3 = D_3 x_{Z/2} Z/12 (ours).
%  - p. 11: « gamma - n beta = alpha » is 2 gamma - n beta = alpha.
%  - p. 16: the third injection repeats the second.
%  - pp. 19, 28, 34: fibre products written over Z/12Z where the maps only
%    exist over Z/2Z (or Z/6Z with Z/12Z); the true forms are given.
%  - p. 22: the invariant ring is polynomial (6 invertible), and
%    54 Delta = gamma_2^3 - 2 gamma_3^2 is the discriminant, not a relation;
%    prod_{i != j}(e_i - e_j) = -Delta; p. 24's « n <= 6 » is n < 6.
%  - p. 29: « rho'_s = rho_f^{-1} » is rho'_s = rho_s^{-1}, as the next line uses.
%  - p. 32: an extension with quotient map alpha.phi_p has class
%    alpha^{-1}.eta_p (immaterial for alpha = -1, the case used).
%  - p. 34: the vertical « = » between Pi~(4) and T~_{11} is an index-2
%    inclusion.
%  - p. 35: the map ST_{11} -> SPi^D(4) has the values the page writes only
%    with the relation rho^3 = sigma^2 = varpi^{-4}.
%
% Announced and never established in the folder: the sigma-invariance of
% Pi~°_n for general n (p. 5); the general formula for sigma~(rho~_i) in the
% extension (p. 5 stops); the order of E's model SL'(2, Z/6Z) (p. 9, « à
% déterminer »); the cotangent guess of p. 19; the axioms for the quadruples
% (A, sigma, A', rho) of p. 31, which refer to pages not in this folder; the
% definition of PM_{11}.
%
% Pass: Opus 5.5 (claude-opus-5-5), 2026-09-24 — first pass, unchecked against the pages by a human.
\documentclass[11pt,a4paper]{article}
% Le préambule commun à toutes les transcriptions du fonds Grothendieck.
%
% Il tient en une idée : **l'incertitude se compose, elle ne se résout pas.**
% Une transcription qui lisse un mot douteux a détruit la seule chose pour
% laquelle on la faisait — un lecteur doit pouvoir savoir, à chaque endroit,
% ce qui était sur le feuillet et ce qui a été ajouté. Les six macros qui
% suivent sont donc l'appareil critique, et non de la décoration.
%
% Les mêmes commandes sont reconnues par `scripts/render.mjs`, qui produit la
% vue de lecture du site. Une macro ajoutée ici doit l'être aussi là-bas,
% faute de quoi le rendu échouera bruyamment — ce qui est voulu : un rendu
% silencieusement faux est pire qu'un rendu absent.

\ProvidesPackage{grothendieck}[2026/08/08 Transcriptions du fonds Grothendieck]

\usepackage{amsmath,amssymb,amsthm}
\usepackage{mathrsfs}
\usepackage{stmaryrd}
% Les diagrammes commutatifs sont partout dans ce fonds : c'est la langue
% même de la géométrie algébrique de Grothendieck, et une transcription qui
% les rendrait en prose ne transcrirait rien.
\usepackage{tikz-cd}
% Français par défaut — c'est la langue des feuillets — mais l'anglais est
% chargé aussi : la traduction et le résumé sont des documents anglais, et la
% typographie française (espaces avant les deux-points) y serait une faute.
% Ces documents basculent par \selectlanguage{english} en tête de corps.
\usepackage[english,french]{babel}
\usepackage{fontspec}
\usepackage{geometry}
\geometry{a4paper,margin=25mm,marginparwidth=28mm}
\usepackage{marginnote}
\usepackage{xcolor}
% ulem, not soul. `soul`'s \st and \ul are built on a character-by-character
% scan that XeLaTeX's font machinery breaks: the document fails with an
% undefined control sequence rather than degrading. ulem does the same two jobs
% and is Unicode-engine safe. `normalem` stops it hijacking \emph, which the
% transcriptions use for Grothendieck's own emphasis.
\usepackage[normalem]{ulem}
\usepackage{hyperref}
\hypersetup{colorlinks=true,linkcolor=black,urlcolor=black,pdfborder={0 0 0}}

% --- Métadonnées du lot -----------------------------------------------------
% Reprises telles quelles par le script de rendu, qui les affiche en tête de
% la vue de lecture. Elles fixent ce que le fichier prétend couvrir : sans
% elles, une transcription flotte sans point d'attache dans un fonds de seize
% mille feuillets.

\newcommand{\folder}[1]{\gdef\grothFolder{#1}}
\newcommand{\batch}[1]{\gdef\grothBatch{#1}}
\newcommand{\leaves}[2]{\gdef\grothFirst{#1}\gdef\grothLast{#2}}
\newcommand{\foldertitle}[1]{\gdef\grothFoldertitle{#1}}
\gdef\grothFolder{?}\gdef\grothBatch{?}\gdef\grothFirst{?}\gdef\grothLast{?}\gdef\grothFoldertitle{}

% --- Le résumé --------------------------------------------------------------
%
% Il ouvre la lecture modernisée, sous le titre « Résumé », et il est composé
% en petit. Ce n'est pas un ornement : c'est le seul passage du document qui
% ne soit pas des mathématiques, et un lecteur doit pouvoir le reconnaître
% avant de l'avoir lu — puis le sauter s'il connaît déjà le sujet, ou s'y
% arrêter s'il ne le connaît pas. Le filet qui le suit marque la reprise du
% corps.

\newenvironment{resume}
  {\small\setlength{\parskip}{0.45em}}
  {\par\medskip\hrule\medskip}

% --- Mots-clés ---------------------------------------------------------------
%
% En anglais, en clôture du résumé de la lecture modernisée : le vocabulaire
% moderne sous lequel chercher ce que les feuillets construisent. Le manifeste
% du site (`npm run manifest`) les extrait pour étiqueter la cote entière —
% cette ligne est la source unique des tags, il n'y a pas de fichier à côté.
\newcommand{\keywords}[1]{\par\smallskip\noindent
  {\footnotesize\textsc{Keywords} --- \textit{#1}}\par}

% --- L'appareil critique ----------------------------------------------------

% Le numéro de feuillet, en marge. C'est l'ancre qui permet de régler un
% désaccord sur une formule en regardant *un* feuillet plutôt qu'une chemise
% de six cents. Le site s'en sert aussi pour tourner les pages du fac-similé
% quand on fait défiler la transcription.
\newcommand{\leaf}[1]{%
  \marginnote{\footnotesize\color{gray!70}\textsf{#1}}%
  \label{leaf:#1}\ignorespaces}

% Illisible : on ne devine pas. Un mot inventé qui se lit comme les autres est
% le pire résultat possible.
\newcommand{\ill}{\textcolor{red!60!black}{[\,\ldots\,]}}

% Lecture douteuse : le mot est donné, et signalé comme incertain.
\newcommand{\uncertain}[1]{\uline{#1}}

% Ajout de l'éditeur — une lettre manquante, un mot sous-entendu.
\newcommand{\add}[1]{\textcolor{blue!50!black}{[#1]}}

% Biffé par Grothendieck lui-même. Ce qu'il a rayé fait partie du document :
% une rature dit souvent où la pensée a changé de direction.
\newcommand{\struck}[1]{\sout{#1}}

% Note du transcripteur, sur l'état matériel du feuillet.
\newcommand{\note}[1]{\par\smallskip{\small\color{gray!60!black}\textit{[#1]}}\par\smallskip}

% Note marginale de Grothendieck, distincte du corps du texte.
\newcommand{\marginal}[1]{\par\smallskip\noindent\hspace{1em}%
  {\small\color{gray!60!black}#1}\par\smallskip}

% --- Titre ------------------------------------------------------------------

\AtBeginDocument{%
  \begin{center}
    {\large\bfseries Fonds Alexandre Grothendieck --- cote n\textsuperscript{o}~\grothFolder}\\[2pt]
    {\itshape\grothFoldertitle}\\[4pt]
    {\small feuillets \grothFirst--\grothLast\ (lot \grothBatch)}\\[2pt]
    {\footnotesize\color{gray!60!black}Transcription établie d'après le fac-similé numérisé
     par l'Université de Montpellier.\\
     Les lectures incertaines sont soulignées, les illisibles notées [\,\ldots\,],
     les ajouts entre crochets.}
  \end{center}
  \vspace{1em}}

\endinput


\folder{140-1}
\pages{1}{36}
\dating{[à partir de 1978-à partir de 1982]}
\watermark{Édition de démonstration\\interprétation personnelle de l'œuvre}
\foldertitle{"Longue Marche" [brouillon] : notes manuscrites (s.d.) — lecture modernisée du dossier entier}

\begin{document}

\section*{Résumé}

\begin{resume}
Ces trente-six pages sont des brouillons : la première porte, de la main de
Grothendieck, le titre « Longue Marche », celui du long manuscrit de 1981, \emph{La
Longue Marche à travers la théorie de Galois}. Elles tournent autour d'un seul
objet, qu'il faut d'abord se représenter sans formules.

Prenez un tore — la surface d'une bouée — et percez-y un petit trou. Les
déformations de cette surface sur elle-même, comptées à déformation continue
près, forment un groupe : le groupe de ses symétries « essentielles ». Ce groupe
est un vieil ami, $SL(2, \mathbb{Z})$, le groupe des matrices entières $2 \times
2$ de déterminant $1$ ; Grothendieck le note $\mathcal{T}_{1,1}$, groupe de
Teichmüller du tore à un trou. Mais si l'on tient fermement le bord du trou — si
l'on interdit de le faire tourner —, un phénomène nouveau apparaît : faire faire
au bord un tour complet, puis le ramener, n'est plus une opération triviale. Comme
une ceinture qu'on vrille en tenant ses deux bouts, la surface garde la mémoire
du nombre de tours. Le groupe s'agrandit d'un compteur de tours, un facteur
$\mathbb{Z}$ qui commute à tout, et devient $S\mathcal{T}_{1,1}$, une
\emph{extension centrale} de $SL(2, \mathbb{Z})$. Ce groupe plus grand est un autre
vieil ami : le groupe des tresses à trois brins.

La question des feuillets est de comprendre ce compteur de tours par un autre
chemin, algébrique celui-là. La sphère privée de trois points, $0$, $1$ et $\infty$,
possède six symétries qui permutent ces trois points ; on passe au quotient, mais
en se souvenant des symétries — comme un kaléidoscope qu'on décrirait par un seul
de ses secteurs et les miroirs qui le bordent. L'objet obtenu, que Grothendieck
note $(U_{0,3}, \mathfrak{S}_3)$ et qu'on appelle aujourd'hui un orbifold ou un
champ, a un groupe fondamental remarquablement simple : deux générateurs $\rho$
et $\sigma$, avec $\rho^3 = \sigma^2 = 1$. C'est encore, déguisé,
$SL(2, \mathbb{Z})$ divisé par $\pm 1$. Les feuillets montrent comment remonter
de ce groupe à $SL(2, \mathbb{Z})$, puis à $S\mathcal{T}_{1,1}$, en ajoutant chaque
fois un peu de « torsion » ; et ils montrent que chacune de ces remontées se lit
sur un nombre : le reste modulo $2$, modulo $4$, modulo $12$ d'un caractère du
groupe. Le nombre $12$ qui revient partout — l'abélianisé de $SL(2, \mathbb{Z})$
est $\mathbb{Z}/12\mathbb{Z}$ — est le fil conducteur des calculs.

Le chemin passe par un détour géométrique. Au-dessus de la sphère à trois trous
vit une droite naturelle, la plus simple des droites tautologiques ; Grothendieck
montre (pages 22 et 24) que, vue avec ses six symétries, elle engendre toutes les
droites possibles et qu'il faut exactement six copies d'elle pour obtenir la droite
triviale. Les points de cette droite privés de zéro forment un espace dont le
groupe fondamental, symétries comprises, est justement le groupe des tresses —
et la dernière page conclut que $S\mathcal{T}_{1,1}$ est bien ce que la
sphère à trois trous et sa droite fournissent.

À côté de ce fil courent deux autres, plus brefs : une page sur les groupes de
symétries d'une surface à bord en général (page 20), et quelques pages sur les
\emph{cartes}, ces graphes dessinés sur une surface que Grothendieck décrit au
même moment par une permutation et une involution d'un ensemble fini — ce sont
les « dessins d'enfants » de l'\emph{Esquisse d'un programme}, et le groupe
$\langle \rho, \sigma \mid \rho^3 = \sigma^2 = 1 \rangle$ y est justement celui
des cartes dont tous les sommets sont de valence trois.

Ce sont des brouillons, et ils en ont les marques : des questions laissées en
suspens (« A-t-on bien $S\mathcal{T}_{11} \simeq S\Pi^{D}_{03}$ ?? »), résolues
quelques feuillets plus loin ; un « à vérifier » barré et remplacé par « OK » ;
une conjecture notée d'un « !! », sur le fibré cotangent, que rien ne vient
confirmer. La datation des archivistes est large ; la seule date écrite sur ces
feuillets, novembre 1981, est d'une autre main.

Les noms à chercher ensuite : groupe modulaire, groupe des tresses, revêtement
universel de $SL(2, \mathbb{R})$, groupe fondamental orbifold, champ des modules des
courbes elliptiques, groupe de Picard équivariant, groupe cartographique, tour de
Teichmüller.

\keywords{central extension, modular group, braid group B3, universal cover of SL(2,R), orbifold fundamental group, quotient stack, moduli stack of elliptic curves, anharmonic group, dicyclic group, mapping class group, boundary Dehn twist, equivariant Picard group, discriminant complement, combinatorial map, cartographic group}
\end{resume}

\pagerange{1}{36}
\section*{Le fil du dossier, et les conventions}

\subsection*{Les stations}

Le dossier est fait de feuillets d'encres et de papiers différents ; une suite
numérotée de sa main (① à ④, « 4 bis », 5, ⑥, ⑦ aux pages 3, 5, 7, 11, 13, 15, 17,
19) en forme l'ossature, et les autres feuillets s'y rattachent. Dans l'ordre où
on les lira :

\begin{itemize}
\item \textbf{Une page isolée} — page 1 : un calcul de matrices $2 \times 2$, sous
le titre « Longue Marche ».
\item \textbf{Les extensions centrales de $\Pi^{\mathbb{D}}_n$} — pages 3 et 5 :
le groupe $\Pi^{\mathbb{D}}_n = \mathbb{Z}/n * \mathbb{Z}/2$, son sous-groupe
libre $\Pi_n$, et ce que deviennent les générateurs de $\Pi_n$ dans l'extension
centrale $\tilde\rho^{\,n} = \tilde\sigma^2$.
\item \textbf{Le cas $n = 3$ et ses quotients finis} — pages 7, 9 et 14 :
$SL(2, \mathbb{Z})$, $GL(2, \mathbb{Z})$, le groupe $\mathbb{D}'_3$ d'ordre $12$,
et des quotients d'ordres $18$, $36$, $72$.
\item \textbf{La normalisation des relèvements} — pages 11, 13 et 18 : pour $n$
impair, l'extension de $\Pi^{\mathbb{D}}_n$ tirée de son abélianisé admet des
générateurs avec $\tilde\rho^{\,n} = \tilde\sigma^2 = \tilde\omega^{\nu}$, et $\nu$
est calculé.
\item \textbf{Les groupes $\mathcal{T}_{1,1}$ et $S\mathcal{T}_{1,1}$} — pages 15,
17 et 19 : $M_{1,1} \simeq (U_{0,3}, \mathbb{D}'_3)$, $\mathcal{T}_{1,1} \simeq
SL(2, \mathbb{Z})$, et $S\mathcal{T}_{1,1}$ comme image inverse de
$\mathbb{Z} \to \mathbb{Z}/12\mathbb{Z}$.
\item \textbf{Surfaces à bord} — page 20 : la suite exacte qui relie
$S\mathcal{T}$, $\mathcal{T}$ et les tours de bord, pour une surface quelconque.
\item \textbf{Cartes} — pages 16, 29 (moitié inférieure) et 31 : sommets, arêtes,
faces, drapeaux ; changement d'orientation ; orbites de $\rho$.
\item \textbf{Le lemme} — pages 22, 24 et 29 (moitié supérieure) :
$H^1(U_{0,3}, \mathfrak{S}_3; \mathbb{G}_m) \simeq \mathbb{Z}/6\mathbb{Z}$, engendré
par $\mathcal{O}(1)$.
\item \textbf{Les champs au-dessus de $M_{1,1}$} — page 26.
\item \textbf{Le diagramme des extensions} — pages 28 et 30 : les quatre groupes
finis $\mathbb{D}_3$, $\mathbb{D}'_3$, $\mathbb{D}''_3$, $\mathbb{D}'''_3$.
\item \textbf{Classes d'extensions, et le caractère du revêtement} — pages 32 et
34.
\item \textbf{La réponse} — pages 33 et 35 : $S\mathcal{T}_{1,1} \simeq
S\Pi^{\mathbb{D}}_{0,3}$, et la place de $S\mathcal{T}_{1,1}$ dans l'extension
$S\Pi^{\mathbb{D}}_{0,3}(4)$.
\end{itemize}

La question posée à la page 19 (« A-t-on bien $S\mathcal{T}_{11} \simeq
S\Pi^{D}_{03}$ ?? ») reçoit sa réponse à la page 33, et la proposition des
pages 11 et 13, faite pour $n$ impair quelconque, est exactement ce qui sert à la
page 34 pour $n = 3$ : c'est le fil qui relie les deux moitiés du dossier. Le
bas de la page 17 porte, tête-bêche et barré, un fragment d'une autre rédaction
(« 8. Groupes et groupoïdes fondamentaux », sur des $\Delta$-complexes
normalisés) ; il est sans rapport avec le reste et n'est pas repris ici.

\subsection*{Les conventions}

\emph{Les groupes $\Pi^{\mathbb{D}}_n$.} Pour $n \geqslant 2$,
\[
\Pi^{\mathbb{D}}_n = \langle \rho, \sigma \mid \rho^n = \sigma^2 = 1 \rangle
\simeq \mathbb{Z}/n\mathbb{Z} * \mathbb{Z}/2\mathbb{Z} ,
\]
et $\mathbb{D}_n$ désigne son quotient diédral, d'ordre $2n$, obtenu en imposant
$\sigma\rho\sigma^{-1} = \rho^{-1}$, c'est-à-dire $(\sigma\rho)^2 = 1$. On pose
$\rho_0 = \sigma\rho\sigma^{-1}\rho = (\sigma\rho)^2$ et
$\rho_i = \rho^i \rho_0 \rho^{-i}$, indices pris modulo $n$. Le noyau $\Pi_n$ de
$\Pi^{\mathbb{D}}_n \to \mathbb{D}_n$ est un groupe libre de rang $n - 1$,
engendré par $\rho_0, \dots, \rho_{n-1}$ avec la seule relation
$\rho_{n-1}\rho_{n-2}\cdots\rho_0 = 1$ : c'est le groupe fondamental
$\Pi_{0,n}$ de la sphère privée de $n$ points.\footnote{Le dossier n'énonce pas
ce fait ; il se lit sur la caractéristique d'Euler : $\Pi^{\mathbb{D}}_n$ est
virtuellement libre de caractéristique $\frac1n + \frac12 - 1$, un sous-groupe
sans torsion d'indice $2n$ est donc libre de rang $1 - 2n(\frac1n + \frac12 - 1)
= n - 1$, et les $n$ éléments $\rho_i$, soumis à une relation, l'engendrent.
Géométriquement : $\mathbb{D}_n$ agit sur la sphère de Riemann par
$z \mapsto \zeta z$, $z \mapsto 1/z$, et $\Pi^{\mathbb{D}}_n$ est le groupe
fondamental orbifold de la sphère privée des racines $n$-ièmes de l'unité,
modulo cette action.} Le cas qui occupe le dossier est $n = 3$ :
\[
\Pi^{\mathbb{D}}_{0,3} := \Pi^{\mathbb{D}}_3, \qquad
\mathbb{D}_3 = \mathfrak{S}_3, \qquad \Pi_{0,3} = \Pi_3 ,
\]
le groupe fondamental de $U_{0,3} = \mathbb{P}^1 - \{0, 1, \infty\}$ et celui de
l'orbifold $(U_{0,3}, \mathfrak{S}_3)$, où $\mathfrak{S}_3$ agit par les six
substitutions anharmoniques de la coordonnée. On a
$\Pi^{\mathbb{D}}_{0,3} \simeq PSL(2, \mathbb{Z})$, isomorphisme que le dossier
utilise sans l'écrire.\footnote{La page 26 écrit $SL(2, \mathbb{Z})'$ pour le
groupe qui correspond à $(U_{0,3}, \mathbb{D}_3)$, et la page 7
$GL(2, \mathbb{Z})'$ pour l'extension de $\mathbb{D}_3 \times \{1, \tau\}$ par
$\Pi_{0,3}$ ; on lit l'accent comme « modulo $\pm 1 $ ». Les pages écrivent
indifféremment $\Pi^{D}$ et $\Pi^{\mathbb{D}}$, $\Pi_{03}$ et $\Pi_{0,3}$,
$\mathcal{T}_{11}$ et $\mathcal{T}_{1,1}$ ; on unifie.}

\emph{Deux extensions centrales à ne pas confondre.} Le dossier donne le même nom
$\widetilde{\Pi}{}^{\mathbb{D}}_n$ à deux groupes distincts. On les sépare :
\begin{itemize}
\item $\widetilde{\Pi}{}^{\mathbb{D}}_n = \langle \tilde\rho, \tilde\sigma \mid
\tilde\rho^{\,n} = \tilde\sigma^2 \rangle$, extension centrale de
$\Pi^{\mathbb{D}}_n$ par le groupe cyclique infini engendré par
$\tilde\omega_0 = \tilde\rho^{\,n} = \tilde\sigma^2$ (pages 3 et 5) ;
\item $S\Pi^{\mathbb{D}}_n = \Pi^{\mathbb{D}}_n \times_{\mathbb{Z}/2n\mathbb{Z}}
\mathbb{Z}$, pour $n$ impair, l'image inverse de l'extension canonique
$0 \to \mathbb{Z} \xrightarrow{2n} \mathbb{Z} \to \mathbb{Z}/2n\mathbb{Z} \to 0$ par
l'abélianisation $\Pi^{\mathbb{D}}_n \to \mathbb{Z}/2n\mathbb{Z}$ (page 11), avec
le générateur $\tilde\omega$ de son noyau.
\end{itemize}
La page 11 appelle le second $\widetilde{\Pi}{}^{D}_n$ ; la page 33 appelle
$S\Pi^{D}_{0,3}$ son cas $n = 3$, et c'est le nom qu'on retient. Pour $n = 3$ les
deux coïncident ; pour $n$ impair plus grand, le premier est un sous-groupe
d'indice $|\nu|$ du second (voir les pages 11 à 13).

\emph{Les $\Pi'$.} À la page 9, $\Pi'_{0,3}$ est le sous-groupe d'indice $3$ de
$\Pi_{0,3}$, noyau du caractère $\rho_0, \rho_1, \rho_\infty \mapsto 1 \in
\mathbb{Z}/3\mathbb{Z}$. À la page 28, le même symbole désigne
$\Pi_{0,3} \times \{1, \dot\omega\}$, c'est-à-dire le sous-groupe de congruence
$\Gamma(2)$ de $SL(2, \mathbb{Z})$ ; on l'écrit $\Gamma(2)$.

\emph{Orbifolds et champs.} Pour un groupe $G$ agissant sur un espace $X$, le
dossier note $(X, G)$ ce qu'on appelle aujourd'hui le champ quotient $[X/G]$ ou
l'orbifold $X/\!/G$, et $\pi_1(X, G; P)$ son groupe fondamental : l'extension de
$G$ par $\pi_1(X, P)$ formée des couples (élément de $G$, chemin de $P$ à son
image). Si $G$ agit à travers un quotient, les éléments du noyau de l'action
restent dans le groupe fondamental : c'est ce qui distingue $(U_{0,3},
\mathbb{D}'_3)$ de $(U_{0,3}, \mathbb{D}_3)$.\footnote{Le mot « orbifold » est
de Thurston (notes de Princeton, 1977--1978) ; les champs algébriques et leur
groupe fondamental sont dans Deligne--Mumford (1969). Grothendieck, dans la
\emph{Longue Marche}, parle de « multiplicités modulaires » ; le dossier ne
nomme ni l'un ni l'autre, et ne donne à $(X, G)$ aucun nom.}

\subsection*{Ce que le dossier annonce et n'établit pas}

\begin{itemize}
\item l'invariance de $\widetilde{\Pi}{}^{\circ}_n$ par $\tilde\sigma$ pour $n$
quelconque (page 5) ; la page calcule deux cas et s'arrête ; la lecture règle le
cas $n$ impair, et le signale ;
\item l'identification du groupe $\mathcal{E}$ d'ordre $36$ à un
« $SL'(2, \mathbb{Z}/6\mathbb{Z})$ » dont l'ordre est « à déterminer » (page 9) ;
\item la suggestion de la page 19, $PM_{1,1} \simeq (T^{*}U_{0,3},
\mathfrak{S}'_3)$, ni démontrée ni reprise, et le symbole $PM_{1,1}$ lui-même,
jamais défini ;
\item les axiomes des quadruples $(A, \sigma, A', \rho)$ de la page 31, qui
renvoient à des pages qui ne sont pas dans ce dossier ;
\item l'équivalence $\mathrm{Pol}_\tau^{\Omega} \simeq \mathrm{Pol}_{\tau/\Omega}$
de la page 16, écrite sans que ses termes soient définis.
\end{itemize}

\pagerange{1}{1}
\section*{Un calcul matriciel (page 1)}

Sous le titre, au crayon, un calcul sur des matrices $2 \times 2$ qui ne se
rattache à rien d'autre du dossier. Soient
\[
\alpha = \begin{pmatrix} a & b \\ c & d \end{pmatrix}, \quad
\delta_\alpha = \det\alpha, \qquad
u = \begin{pmatrix} \mu & \lambda \\ 0 & 1 \end{pmatrix}, \qquad
\rho = \begin{pmatrix} 0 & v_1 \\ v & w \end{pmatrix},
\]
et $\underline{a} = \alpha^{-1}u$. Pour toute matrice $2 \times 2$ et tous
scalaires $x, y$,
\[
\det(x\rho + y) = x^2 \det\rho + xy\,\mathrm{Tr}\,\rho + y^2
= -v v_1 x^2 + w xy + y^2 ,
\]
ce qui, pour $v_1 = 1$, $x = d$, $y = c$, est l'expression $-vd^2 + wcd + c^2$
écrite en tête de la page. On calcule
\[
\alpha^{-1}u = \frac{1}{\delta_\alpha}
\begin{pmatrix} d\mu & d\lambda - b \\ -c\mu & a - c\lambda \end{pmatrix} .
\]
\emph{Si $\underline{a}$ commute à $\rho$} et si $v \neq 0$, alors $\underline{a}$
est un polynôme en $\rho$ — le commutant d'une matrice $2 \times 2$ non scalaire
est la sous-algèbre qu'elle engendre —, soit $\underline{a} = x\rho + y$ ; la
première colonne de $x\rho + y$ étant $(y, xv)$, on identifie
\[
\underline{a} = \frac{\mu}{\delta_\alpha}\Bigl(-\frac{c}{v}\,\rho + d\Bigr) ,
\]
qui est la formule encadrée.\footnote{La page écrit la formule sans hypothèse ;
elle identifie les coefficients sur la première colonne seulement. Les deux
autres coefficients imposent $d\lambda - b = -c\mu v_1/v$ et
$a - c\lambda = \mu(d - cw/v)$, ce qui est exactement la condition de
commutation. Les étapes intermédiaires de la page nomment tantôt $c_0, d_0$ les
coefficients de $\rho$ et de $1$, tantôt l'inverse ; une ligne donne pour $c_0$
une fraction de dénominateur « $-v^2 + wcd + c^2$ », premier terme de lecture
incertaine, qui n'est pas utilisée ensuite.} Son déterminant se calcule de deux
façons :
\[
\det \underline{a} = \frac{\mu}{\delta_\alpha}
\qquad\text{et}\qquad
\det \underline{a} = \Bigl(\frac{\mu}{\delta_\alpha}\Bigr)^{2}
\Bigl[-\frac{c^2}{v^2}\,v v_1 - \frac{cd}{v}\,w + d^2\Bigr] ,
\]
la première directement, la seconde par la formule du déterminant ; leur égalité
est une relation entre les données.\footnote{La page pose la première en question
(« $\det a \overset{?}{=} \delta_\alpha^{-1}\mu$ ») : elle est vraie sans
hypothèse, $\det u = \mu$. Dans la seconde, la page écrit $+\frac{dc}{v}w$, le
signe étant un $+$ repassé ; avec $x = -c/v$ et $y = d$ le terme $xyw$ vaut
$-\frac{cd}{v}w$. La ligne s'arrête sur des points de suspension.} La page ne dit
pas à quoi sert ce calcul.

\pagerange{3}{5}
\section*{Les extensions centrales de $\Pi^{\mathbb{D}}_n$ (pages 3 et 5)}

\subsection*{Les relèvements des $\rho_i$ (page 3)}

Dans $\widetilde{\Pi}{}^{\mathbb{D}}_n = \langle \tilde\rho, \tilde\sigma \mid
\tilde\rho^{\,n} = \tilde\sigma^2 \rangle$, l'élément
$\tilde\omega_0 = \tilde\rho^{\,n} = \tilde\sigma^2$ commute aux deux générateurs,
donc est central, et engendre un sous-groupe cyclique infini : on a une extension
centrale
\[
1 \to \mathbb{Z}\tilde\omega_0 \to \widetilde{\Pi}{}^{\mathbb{D}}_n \to
\Pi^{\mathbb{D}}_n \to 1 .
\]
On relève les $\rho_i$ par
\[
\tilde\rho_0 = \tilde\sigma\tilde\rho\tilde\sigma^{-1}\tilde\rho
= (\tilde\sigma\tilde\rho)^2\,\tilde\omega_0^{-1}, \qquad
\tilde\rho_i = \tilde\rho^{\,i}\,\tilde\rho_0\,\tilde\rho^{-i} ,
\]
la seconde égalité parce que $\tilde\sigma^{-1} = \tilde\sigma\,\tilde\omega_0^{-1}$.
Comme $\tilde\rho^{\,n}$ est central, $\tilde\rho_{i+n} = \tilde\rho_i$ : les
indices vivent modulo $n$. Le produit des $\tilde\rho_i$ se calcule en
télescopant :
\[
\tilde\rho_{n-1}\tilde\rho_{n-2}\cdots\tilde\rho_0
= \tilde\rho^{\,n-1}(\tilde\rho_0\tilde\rho^{-1})^{n-1}\tilde\rho_0
= \tilde\rho^{\,n}\,(\tilde\rho^{-1}\tilde\rho_0)^n .
\]
Or $\tilde\rho^{-1}\tilde\rho_0 = g\,\tilde\rho\,g^{-1}$ avec
$g = \tilde\rho^{-1}\tilde\sigma$, donc $(\tilde\rho^{-1}\tilde\rho_0)^n =
g\,\tilde\rho^{\,n}g^{-1} = \tilde\omega_0$, et
\[
\boxed{\tilde\rho_{n-1}\tilde\rho_{n-2}\cdots\tilde\rho_0 = \tilde\omega_0^{\,2}} .
\]
Là où le produit des $\rho_i$ vaut $1$ en bas, celui des relèvements vaut
$\tilde\omega_0^2$, non $\tilde\omega_0$ : c'est le premier des facteurs $2$ du
dossier.

Soit $\widetilde{\Pi}_n$ l'image inverse de $\Pi_n$ dans
$\widetilde{\Pi}{}^{\mathbb{D}}_n$.\footnote{La page commence par le définir comme
le sous-groupe engendré par les $\tilde\rho_i$, biffe, et le redéfinit comme image
inverse : les deux ne coïncident pas, le premier est d'indice $2$ dans le second,
et c'est l'objet de la suite.} C'est une extension centrale de $\Pi_n$ par
$\mathbb{Z}\tilde\omega_0$, présentée par
\[
\widetilde{\Pi}_n = \langle \tilde\rho_0, \dots, \tilde\rho_{n-1}, \tilde\omega_0
\mid \tilde\omega_0 \text{ central},\
\tilde\rho_{n-1}\cdots\tilde\rho_0 = \tilde\omega_0^{\,2} \rangle .
\]
Soit de même
\[
\widetilde{\Pi}{}^{\circ}_n = \langle \tilde\rho_0, \dots, \tilde\rho_{n-1} \mid
\varpi = \tilde\rho_{n-1}\cdots\tilde\rho_0 \text{ central} \rangle ,
\]
la condition revenant à demander que les permutations circulaires du produit
soient égales. C'est une extension centrale de $\Pi_n$ par $\mathbb{Z}\varpi$, et
\[
\widetilde{\Pi}_n \simeq \widetilde{\Pi}{}^{\circ}_n \wedge^{\mathbb{Z}\varpi}
\mathbb{Z}\tilde\omega_0, \qquad \varpi \longmapsto \tilde\omega_0^{\,2} ,
\]
la somme amalgamée le long de $\mathbb{Z} \xrightarrow{2} \mathbb{Z}$ :
$\widetilde{\Pi}_n$ se déduit de $\widetilde{\Pi}{}^{\circ}_n$ en adjoignant une
racine carrée centrale de $\varpi$. En particulier $\widetilde{\Pi}{}^{\circ}_n$ est
le sous-groupe de $\widetilde{\Pi}_n$ engendré par les $\tilde\rho_i$, et il y est
d'indice $2$.\footnote{Comme $\Pi_n$ est libre, toutes ces extensions sont
scindées en tant qu'extensions abstraites : $\widetilde{\Pi}_n \simeq \Pi_n \times
\mathbb{Z}$. Ce qui porte l'information n'est donc pas la classe de l'extension
mais la position des relèvements canoniques $\tilde\rho_i$ des générateurs
géométriques, dont le produit vaut $\tilde\omega_0^2$ et non $1$. La page note
$\wedge_{\mathbb{Z}\varpi}$ la somme amalgamée ; au-dessus de $\Pi_n$ dans la
suite exacte, un mot lu « can ».}

\subsection*{L'action de $\sigma$ (page 5)}

La question de la page est : le sous-groupe $\widetilde{\Pi}{}^{\circ}_n$, distingué
dans $\widetilde{\Pi}_n$ et stable par conjugaison par $\tilde\rho$, l'est-il par
$\tilde\sigma$, donc dans $\widetilde{\Pi}{}^{\mathbb{D}}_n$ tout entier ? En bas,
dans $\Pi^{\mathbb{D}}_n$, on a $\sigma\rho\sigma^{-1} = \rho_0\rho^{-1}$, d'où
\[
\sigma(\rho_0) = \rho_1, \quad \sigma(\rho_1) = \rho_0, \quad
\sigma(\rho_2) = \rho_0\,\rho_{n-1}\,\rho_0^{-1}, \quad \dots,
\]
\[
\sigma(\rho_i) = (\rho_n\rho_{n-1}\cdots\rho_{n+2-i})\,\rho_{n+1-i}\,
(\rho_n\rho_{n-1}\cdots\rho_{n+2-i})^{-1} ,
\]
où $\sigma(x) = \sigma x \sigma^{-1}$ et les indices sont pris modulo $n$ ;
$\sigma$ étant une involution, $\sigma = \sigma^{-1}$ en tant qu'automorphisme
intérieur.\footnote{La formule générale est celle de la page, où les indices
sont surchargés ; elle se vérifie en écrivant
$(\rho_0\rho^{-1})^i = \rho_0\rho_{-1}\cdots\rho_{1-i}\,\rho^{-i}$.} Dans
l'extension, les mêmes formules valent à un facteur central près :
$\tilde\sigma(\tilde\rho_i) = (\cdots)\,\tilde z_i$ avec
$\tilde z_i = \tilde\omega_0^{\,\nu_i}$. L'abélianisé de
$\widetilde{\Pi}{}^{\mathbb{D}}_n$ contient $\mathbb{Z}\tilde\omega_0$ et
$\tilde\sigma$ fixe $\tilde\omega_0^2 = \prod\tilde\rho_i$, d'où
$\sum_i \nu_i = 0$. Comme $\tilde\omega_0 \notin \widetilde{\Pi}{}^{\circ}_n$,
\[
\widetilde{\Pi}{}^{\circ}_n \text{ est stable par } \tilde\sigma
\iff \nu_i \equiv 0 \pmod 2 \text{ pour tout } i .
\]
La page calcule les deux premiers : en utilisant $\tilde\sigma^{\pm 2} =
\tilde\omega_0^{\pm1} = \tilde\rho^{\,\pm n}$,
\[
\tilde\rho_1^{-1}\,\tilde\sigma(\tilde\rho_0)
= \tilde\sigma\tilde\rho^{-1}\tilde\sigma^{-2}\tilde\rho\tilde\sigma
= \tilde\omega_0^{-1}\tilde\sigma^{2} = 1,
\qquad
\tilde\rho_0^{-1}\,\tilde\sigma^{-1}(\tilde\rho_1) = 1 ,
\]
donc $\nu_0 = \nu_1 = 0$, puis s'arrête sur le cas général.\footnote{La page
annonce qu'on calculera « au moins $z_0$, $z_1$ », la phrase étant en partie
illisible ; elle écrit ensuite $\tilde\rho_i = \tilde\rho^{\,i}\tilde\sigma
\tilde\rho\tilde\sigma^{-1}\tilde\rho^{\,1-i}$ et $\tilde\sigma(\tilde\rho_i)$, et
s'interrompt. Dans le second calcul la lecture d'un facteur ($\tilde\sigma$ ou
$\tilde\sigma^{-1}$) est incertaine ; le résultat $1$ est juste.}

Pour $n$ impair, la réponse est oui, et elle se lit sur l'abélianisé.\footnote{Cet
argument est de la présente lecture ; la page ne l'a pas.} Si $n$ est impair,
$(\widetilde{\Pi}{}^{\mathbb{D}}_n)_{\mathrm{ab}} = \mathbb{Z}^2/(n, -2) \simeq
\mathbb{Z}$, par $\tilde\rho \mapsto 2$, $\tilde\sigma \mapsto n$ ; alors
$\tilde\rho_i \mapsto 4$ et $\tilde\omega_0 \mapsto 2n \equiv 2 \pmod 4$. Le
caractère $\widetilde{\Pi}{}^{\mathbb{D}}_n \to \mathbb{Z}/4\mathbb{Z}$ qui s'en
déduit, restreint à $\widetilde{\Pi}_n$, tue les $\tilde\rho_i$ et non
$\tilde\omega_0$ : son noyau dans $\widetilde{\Pi}_n$ est donc
$\widetilde{\Pi}{}^{\circ}_n$, qui est ainsi l'intersection de $\widetilde{\Pi}_n$ avec
un sous-groupe distingué de $\widetilde{\Pi}{}^{\mathbb{D}}_n$, et distingué
lui-même. Pour $n$ pair, aucun caractère ne sépare ainsi $\tilde\omega_0$ des
$\tilde\rho_i$, puisque $\tilde\omega_0 = \tilde\rho^{\,n}$ a alors même image
abélianisée que $\tilde\rho_0^{\,n/2}$, et la question reste ouverte dans le
dossier.

\pagerange{7}{9}
\section*{Le cas $n = 3$ : $SL(2, \mathbb{Z})$, $\mathbb{D}'_3$ et leurs quotients finis (pages 7, 9 et 14)}

\subsection*{La tour de la page 7}

Pour $n = 3$, le groupe $\widetilde{\Pi}{}^{\mathbb{D}}_3 = \langle \tilde\rho,
\tilde\sigma \mid \tilde\rho^{\,3} = \tilde\sigma^2 \rangle$ est l'image inverse
$\widetilde{SL}(2, \mathbb{Z})$ de $SL(2, \mathbb{Z})$ dans le revêtement universel
de $SL(2, \mathbb{R})$ (page 17), et $\widetilde{SL}(2, \mathbb{Z}) /
\langle \omega^2 \rangle = SL(2, \mathbb{Z})$, où $\omega = \tilde\omega_0$. La page
organise ses sous-groupes distingués en une tour, chaque flèche étiquetée par le
quotient :
\[
\mathbb{Z}\omega^2 \subset \widetilde{\Pi}{}^{\circ}_{0,3} \subset
\widetilde{\Pi}_{0,3} \subset \widetilde{SL}(2, \mathbb{Z}) \subset
\widetilde{GL}(2, \mathbb{Z}),
\qquad
\text{quotients successifs } \Pi_{0,3},\ \mathbb{Z}/2,\ \mathbb{D}_3,\ \mathbb{Z}/2 ,
\]
avec aussi $\mathbb{Z}\omega^2 \subset \mathbb{Z}\omega \subset
\widetilde{\Pi}_{0,3}$ (quotients $\mathbb{Z}/2$, puis $\Pi_{0,3}$), et
$\widetilde{GL}(2, \mathbb{Z}) \supset \widetilde{\Pi}{}^{\tau}_{0,3}$ de quotient
$\mathbb{D}_3$. Modulo $\omega^2$ : $\widetilde{SL}(2, \mathbb{Z})/\mathbb{Z}\omega^2 =
SL(2, \mathbb{Z})$, $\widetilde{GL}(2, \mathbb{Z})/\mathbb{Z}\omega^2 =
GL(2, \mathbb{Z})$, et
\[
\widetilde{SL}(2, \mathbb{Z}) / \widetilde{\Pi}{}^{\circ}_{0,3} \simeq \mathbb{D}'_3 ,
\]
le groupe d'ordre $12$
\[
\mathbb{D}'_3 = \langle \rho, \sigma \mid \rho^3 = \sigma^2,\ \sigma^4 = 1,\
\sigma\rho\sigma^{-1} = \rho^{-1} \rangle ,
\qquad 1 \to \mathbb{Z}/2 \to \mathbb{D}'_3 \to \mathbb{D}_3 \to 1 ,
\]
le noyau étant engendré par $\omega = \rho^3 = \sigma^2$. C'est le groupe
dicyclique d'ordre $12$, $\mathbb{Z}/3 \rtimes \mathbb{Z}/4$, et c'est aussi le
produit fibré $\mathbb{D}_3 \times_{\mathbb{Z}/2} \mathbb{Z}/4$ de la page 15.
L'extension n'est pas scindée : les deux relèvements de $\sigma$, à savoir $\sigma$
et $\sigma\omega = \sigma^{-1}$, sont d'ordre $4$, et aucun ne peut être l'image
d'une section. Elle diffère en cela de $\mathbb{D}_6 \simeq \mathbb{D}_3 \times
\mathbb{Z}/2 \to \mathbb{D}_3$, qui l'est.\footnote{La note marginale de la page
dit que l'extension n'est pas scindée parce que « les $\rho_{\mathbb{D}'_3}$ qui
relèvent $\rho_{\mathbb{D}_3}$ sont d'ordre 6, non 3 ». C'est inexact : $\rho$
est d'ordre $6$, mais l'autre relèvement $\rho\omega = \rho^4$ est d'ordre $3$.
La conclusion est juste, et la raison est du côté de $\sigma$.}

Au niveau de $SL(2, \mathbb{Z})$ et $GL(2, \mathbb{Z})$, la page écrit
\[
1 \to \mathbb{Z}/2 \times \Pi_{0,3} \to SL(2, \mathbb{Z}) \to \mathbb{D}_3 \to 1,
\qquad
1 \to \mathbb{Z}/2 \times \Pi_{0,3} \to GL(2, \mathbb{Z}) \to \mathbb{D}_3 \times
\{1, \tau\} \to 1 ,
\]
la première étant la réduction modulo $2$, de noyau $\Gamma(2) = \{\pm 1\} \times
\Pi_{0,3}$, la seconde $g \mapsto (g \bmod 2, \det g)$, avec $\tau$ le
déterminant $-1$. Puis, en retirant le facteur $\{\pm 1\}$ :
\[
1 \to \Pi_{0,3} \to GL(2, \mathbb{Z}) \to \mathbb{D}'_3 \rtimes \{1, \tau\} \to 1,
\qquad
1 \to \Pi_{0,3} \to PGL(2, \mathbb{Z}) \to \mathbb{D}_3 \times \{1, \tau\} \to 1 ,
\]
et $PGL(2, \mathbb{Z}) \simeq \pi_1(U_{0,3}, \mathbb{D}_3 \times \{1, \tau\}; P_+)$,
où $\tau$ agit sur $U_{0,3}$ par la conjugaison complexe, qui commute aux
substitutions anharmoniques.\footnote{La page écrit
$\mathbb{D}'_3 \times \{1, \tau\}$ dans la première suite. Ce ne peut être un
produit direct : un élément de déterminant $-1$, tel que
$\mathrm{diag}(1, -1)$, agit sur l'abélianisé $\mathbb{Z}/12$ de
$SL(2, \mathbb{Z})$ par $-1$, donc sur l'abélianisé $\mathbb{Z}/4$ de
$\mathbb{D}'_3$ par $-1$, ce qui n'est pas un automorphisme intérieur. Le quotient
est un produit semi-direct, scindé par $\mathrm{diag}(1, -1)$. La page note
$GL(2, \mathbb{Z})'$ ce que l'on écrit $PGL(2, \mathbb{Z})$ ; la lecture de
$U_{0,3}$ dans $\pi_1(U_{0,3}, \ldots)$ est incertaine. Le point base $P_+$ est
lu ici comme l'un des deux points de $U_{0,3}$ fixés par les éléments d'ordre
$3$ de $\mathfrak{S}_3$, les racines primitives sixièmes de l'unité ; la page ne
le définit pas, mais la page 15 le place sous le point $e^{2i\pi/6}$ du
demi-plan de Poincaré, ce qui s'accorde avec cette lecture.}

\subsection*{Les quotients d'ordres 18 et 36 (page 9)}

La page 9 est une page de schémas ; on en donne les énoncés. Soit
$\Pi'_{0,3} \subset \Pi_{0,3}$ le noyau du caractère $\rho_0, \rho_1, \rho_\infty
\mapsto 1 \in \mathbb{Z}/3$. Il est stable par les substitutions anharmoniques,
qui permutent les $\rho_i$ à conjugaison près, donc distingué dans
$\Pi^{\mathbb{D}}_{0,3}$, d'indice $18$ :
\[
\Pi^{\mathbb{D}}_{0,3}/\Pi'_{0,3} = \langle \rho, \sigma \mid
\rho^3 = \sigma^2 = (\sigma\rho)^6 = 1,\ (\sigma\rho)^2 = (\rho\sigma)^2 \rangle
\xrightarrow{\ \sim\ }
\mathbb{Z}/6 \times_{\mathbb{Z}/2} \mathbb{D}_3
\simeq (\mathbb{F}_3 \times \mathbb{F}_3) \rtimes \{1, \sigma\} ,
\]
par $\rho \mapsto (4, \rho)$, $\sigma \mapsto (3, \sigma)$ ; la relation
$(\sigma\rho)^2 = (\rho\sigma)^2$ dit $\rho_0 = \rho_1$, et $(\sigma\rho)^6 = 1$
dit $\rho_0^3 = 1$. Le produit fibré est pris au-dessus des deux surjections
naturelles sur $\mathbb{Z}/2$, la réduction et la signature ; son sous-groupe
d'indice $2$ est $\mathbb{Z}/3 \times \mathfrak{A}_3$, et $(3, \sigma)$ y agit
trivialement sur le premier facteur et par inversion sur le second.

Au-dessus, dans $SL(2, \mathbb{Z})$, le quotient par le même $\Pi'_{0,3}$ est un
groupe $\mathcal{E}$ d'ordre $36$ :
\[
\mathcal{E} = \langle \rho, \sigma \mid \rho^3 = \sigma^2 = \omega,\ \omega^2 = 1,\
\rho_0 = (\sigma\rho)^2\omega \text{ fixé par } \rho \text{ et d'ordre } 3 \rangle,
\]
\[
\mathcal{E}/\{1, \omega\} \simeq \Pi^{\mathbb{D}}_{0,3}/\Pi'_{0,3} ,
\]
avec, vers l'abélianisé $SL(2, \mathbb{Z})_{\mathrm{ab}} \simeq \mathbb{Z}/12$,
\[
\rho \mapsto 10, \quad \sigma \mapsto 3, \quad \omega \mapsto 6, \quad
\varepsilon_0 = \sigma\rho \mapsto 1, \quad \rho_0 \mapsto 8 .
\]
La page propose d'identifier $\mathcal{E}$ à un sous-groupe
« $SL'(2, \mathbb{Z}/6\mathbb{Z})$ » dont l'ordre reste « à
déterminer ».\footnote{La page écrit aussi
« $SL(2, \mathbb{Z}) \to SL(2, \mathbb{Z})/(\pm 1) \simeq \mathcal{E}$ », de
lecture incertaine ; ce qui est vrai est que $\mathcal{E}$ modulo $\pm 1$ est le
groupe d'ordre $18$. Elle donne pour $\mathbb{D}'_3$, au même endroit, la
présentation $\langle \rho, \sigma \mid \rho^3 = \sigma^2, \rho^6 = 1 \rangle$,
qui est celle de $SL(2, \mathbb{Z})$ ; il manque la relation
$\sigma\rho\sigma^{-1} = \rho^{-1}$ de la page 7.} $\mathcal{E}$ ne peut pas être
un quotient de $SL(2, \mathbb{Z}/6)$ par la réduction : son abélianisé est
$\mathbb{Z}/12$, celui de $SL(2, \mathbb{Z}/6) = SL(2, \mathbb{F}_2) \times
SL(2, \mathbb{F}_3)$ est $\mathbb{Z}/6$. Ce qui est vrai est que
\[
\mathcal{E} \xrightarrow{\ \sim\ } \mathbb{D}_3 \times_{\mathbb{Z}/2}
\mathbb{Z}/12, \qquad g \longmapsto (g \bmod 2,\ g_{\mathrm{ab}}) ,
\]
le groupe que la page 28 appellera $\mathbb{D}'''_3$ : l'application
$SL(2, \mathbb{Z}) \to \mathbb{D}_3 \times \mathbb{Z}/12$ a pour noyau
$\Gamma(2) \cap \ker(\mathrm{ab})$, qui est $\Pi'_{0,3}$ puisque $-1 \mapsto 6$
et que l'abélianisation restreinte à $\Pi_{0,3}$ est le caractère
$\rho_i \mapsto 8$, et son image est d'ordre $36$.\footnote{Ce rapprochement
entre la page 9 et la page 28 est de la présente lecture ; le dossier ne le fait
pas.}

\pagerange{14}{14}
\subsection*{Les mêmes quotients, repris (page 14)}

La page 14 reprend ces identifications. Le quotient d'ordre $18$ y est décrit
comme
\[
\mathbb{Z}/6 \times_{\mathbb{Z}/2} \mathbb{D}_3 = \{ (\nu, g) \in \mathbb{Z}/6
\times \mathbb{D}_3 \mid \mathrm{sg}(g) \equiv \nu \pmod 2 \},
\]
$\mathrm{sg}$ la signature. L'abélianisé $(\Pi_{0,3})_{\mathrm{ab}} =
\mathbb{Z}^3/\mathbb{Z}(1,1,1)$ sur les classes de $\rho_0, \rho_1, \rho_\infty$
donne, tensorisé par $\mathbb{F}_3$, l'espace $\mathbb{F}_3^3/\mathbb{F}_3$ où se
lit le caractère qui définit $\Pi'_{0,3}$, et $\Pi'_{0,3}$ est engendré, comme
sous-groupe distingué, par $\rho_0\rho_1^{-1}$, $\rho_1\rho_\infty^{-1}$,
$\rho_\infty\rho_0^{-1}$ et $\rho_0^3$. Enfin
\[
SL(2, \mathbb{Z}/6)/\{\pm 1\} \simeq \bigl(SL(2, \mathbb{Z}/2) \times
SL(2, \mathbb{Z}/3)\bigr)/\{\pm 1\} ,
\]
d'ordre $6 \cdot 24/2 = 72$, et l'on retrouve, dans $SL(2, \mathbb{Z})$,
$\rho_0 = \sigma\rho\sigma^{-1}\rho = (\sigma\rho)^2\omega$.\footnote{La page porte
encore des produits isolés, en partie illisibles, et un commutateur
$[\rho_0, \rho_1]$ développé puis biffé.}

\pagerange{11}{13}
\section*{Normaliser les relèvements : $\tilde\rho^{\,n} = \tilde\sigma^2 = \tilde\omega^{\nu}$ (pages 11, 13 et 18)}

\subsection*{L'énoncé (page 11)}

Soit $n$ impair. L'abélianisé de $\Pi^{\mathbb{D}}_n$ est $\mathbb{Z}/n \times
\mathbb{Z}/2 \simeq \mathbb{Z}/2n$, et l'on choisit l'isomorphisme
\[
\Pi^{\mathbb{D}}_n \longrightarrow \mathbb{Z}/2n, \qquad
\rho \longmapsto n + 1 = 1 - n, \quad \sigma \longmapsto n .
\]
Par image inverse de $0 \to \mathbb{Z} \xrightarrow{2n} \mathbb{Z} \to \mathbb{Z}/2n
\to 0$, on obtient l'extension centrale
\[
1 \to \mathbb{Z}\tilde\omega \to S\Pi^{\mathbb{D}}_n = \Pi^{\mathbb{D}}_n
\times_{\mathbb{Z}/2n} \mathbb{Z} \to \Pi^{\mathbb{D}}_n \to 1, \qquad
\tilde\omega = (1, 2n) .
\]

\textbf{Proposition} ($n$ impair).\footnote{L'énoncé de la page est écrit en diagonale et une grande partie des mots
de liaison est illisible ; on le reconstitue d'après les calculs des pages 11 et
13, qui le démontrent. Les deux valeurs « $-\frac{n-1}{2}$, $\frac{n+1}{2}$ » et
« $\nu$ impair » sont lisibles ; « $\nu_0$ unique » et « uniquement déterminés »
aussi. La page écrit, au-dessus, une suite
$\mathbb{Z}\tilde\omega^2 \to \widetilde{\Pi}{}^{D}_n \to \Pi^{D}_n \to 1$, dont
le premier terme est vraisemblablement $\mathbb{Z}\tilde\omega$.} \emph{Il existe des relèvements $\tilde\rho$,
$\tilde\sigma$ de $\rho$, $\sigma$ dans $S\Pi^{\mathbb{D}}_n$ tels que
$\tilde\rho^{\,n} = \tilde\sigma^2$. Pour un tel couple,
$\tilde\rho^{\,n} = \tilde\sigma^2 = \tilde\omega^{\nu}$, où la classe de $\nu$
modulo $2n$ ne dépend pas du couple, et vaut $\frac{n^2+1}{2}$ ; les couples
possibles sont les $(\tilde\rho\,\tilde\omega^{2x}, \tilde\sigma\,\tilde\omega^{nx})$,
$x \in \mathbb{Z}$, qui changent $\nu$ en $\nu + 2nx$. On peut donc choisir
$\nu \in \{-\frac{n-1}{2}, \frac{n+1}{2}\}$, et c'est l'unique élément impair de
cet ensemble ; les relèvements sont alors uniquement déterminés, et}
\[
S\Pi^{\mathbb{D}}_n = \langle \tilde\rho, \tilde\sigma, \tilde\omega \mid
\tilde\omega \text{ central},\ \tilde\rho^{\,n} = \tilde\sigma^2 = \tilde\omega^{\nu}
\rangle .
\]


\emph{Démonstration} (pages 11 et 13). Posons $\tilde\rho_0 = (\rho, n+1)$ et
$\tilde\sigma_0 = (\sigma, n)$. Alors $\tilde\sigma_0^{\,2} = (1, 2n) =
\tilde\omega$ et $\tilde\rho_0^{\,n} = (1, n(n+1)) = \tilde\omega^{(n+1)/2}$, de
sorte que $\tilde\rho_0^{\,n}\tilde\sigma_0^{-2} = \tilde\omega^{\alpha}$ avec
$\alpha = \frac{n-1}{2}$. On cherche $\tilde\rho = \tilde\rho_0\tilde\omega^{\beta}$,
$\tilde\sigma = \tilde\sigma_0\tilde\omega^{\gamma}$ avec
$\tilde\rho^{\,n} = \tilde\sigma^2$, c'est-à-dire\footnote{La page écrit « $\gamma - n\beta = \alpha$ » ; la ligne précédente,
$\tilde\rho_0^{\,n}\tilde\sigma_0^{-2} = \tilde\omega^{2\gamma - n\beta}$, et la
solution qu'elle trouve demandent $2\gamma$. De même l'accolade
« $\underbrace{\alpha(\alpha+1)\cdot 2}_{\gamma}$ » désigne le terme dont
$\gamma = \alpha(\alpha+1)$ est le coefficient.}
\[
2\gamma - n\beta = \alpha .
\]
 Comme $n = 2\alpha + 1$, on a
$2(\alpha + 1) - n = 1$, d'où la solution $\beta = \alpha$,
$\gamma = \alpha(\alpha+1) = \frac{n^2-1}{4}$, et
\[
\tilde\rho = (\rho, n^2 + 1), \qquad
\tilde\sigma = \Bigl(\sigma, \frac{n^3+n}{2}\Bigr), \qquad
\tilde\rho^{\,n} = \tilde\sigma^2 = (1, n^3 + n) = \tilde\omega^{\frac{n^2+1}{2}} .
\]
Deux couples solutions diffèrent par $(\tilde\omega^{\beta}, \tilde\omega^{\gamma})$
avec $n\beta = 2\gamma$, soit, $n$ étant impair, $\beta = 2x$ et $\gamma = nx$ ;
l'exposant devient alors $\nu + 2nx$. Reste à réduire $\frac{n^2+1}{2} =
2\alpha^2 + 2\alpha + 1$ modulo $2n = 4\alpha + 2$ :
\[
\frac{n^2+1}{2} \equiv
\begin{cases}
\alpha + 1 = \frac{n+1}{2} & \text{si } \alpha \text{ est pair}, \\[2pt]
-\alpha = -\frac{n-1}{2} & \text{si } \alpha \text{ est impair},
\end{cases}
\pmod{2n}
\]
et dans les deux cas l'exposant obtenu est impair. Les deux éléments de
$\{-\frac{n-1}{2}, \frac{n+1}{2}\}$ diffèrent de $n$, impair : un seul est impair.
La présentation s'en déduit, $\tilde\rho$ et $\tilde\sigma$ relevant les
générateurs de $\Pi^{\mathbb{D}}_n$ et $\tilde\omega$ engendrant le
noyau.\footnote{Les exemples écrits en marge de la page 13 sont les cas $n = 3$
($\frac{n^2+1}{2} = 5 \equiv -1 \bmod 6$) et $n = 5$ ($13 \equiv 3 \bmod 10$).}

Deux conséquences, que la page ne tire pas.\footnote{Elles sont de la présente
lecture.} D'abord, l'application $\widetilde{\Pi}{}^{\mathbb{D}}_n \to
S\Pi^{\mathbb{D}}_n$ qui envoie les générateurs de la page 3 sur les relèvements
normalisés envoie $\tilde\omega_0$ sur $\tilde\omega^{\nu}$ ; elle est injective,
d'image d'indice $|\nu|$. Pour $n = 3$, $\nu = -1$ et les deux extensions
coïncident : c'est la raison pour laquelle $SL(2, \mathbb{Z})$ et son revêtement
$\widetilde{SL}(2, \mathbb{Z})$ se laissent décrire à la fois par la relation
$\tilde\rho^{\,3} = \tilde\sigma^2$ et par l'abélianisé ; pour $n = 5$, $\nu = 3$
et elles diffèrent. Ensuite, pour $n = 3$, les relèvements normalisés sont
\[
\tilde\rho = (\rho, -2), \qquad \tilde\sigma = (\sigma, -3), \qquad
\tilde\rho^{\,3} = \tilde\sigma^2 = \tilde\omega^{-1} ,
\]
dans $\Pi^{\mathbb{D}}_{0,3} \times_{\mathbb{Z}/6} \mathbb{Z}$ : ce sont exactement
ceux de la page 34.

\pagerange{18}{18}
\subsection*{Une reprise (page 18)}

La page 18, au crayon, refait la normalisation en partant de relèvements
satisfaisant $\tilde\sigma^2 = \tilde\rho^{\,n}$ et en essayant
$\tilde\rho \mapsto \tilde\rho\,\tilde\omega^{-1}$, $\tilde\sigma \mapsto
\tilde\sigma\,\tilde\omega^{s}$ ; elle aboutit à la condition
« $1 + 2s = \frac{1-n}{2}$, OK si $\frac{1-n}{2}$ impair », qui est le second cas
de la page 13.\footnote{La page mêle $\tilde\omega$ et $\omega$ d'une ligne à
l'autre ($\tilde\rho_0^{\,n} = \tilde\omega^{1-n} = \omega^{\frac{1-n}{2}}$,
comme si $\omega = \tilde\omega^2$), puis compare un exposant de $\tilde\omega$ à
un exposant de $\omega$ ; la comptabilité n'est pas cohérente telle qu'elle est
écrite, et la lecture ne la répare pas. La première ligne,
« $\tilde\rho_0^{\,n} = \tilde\sigma_0^{\,n} =$ », reste inachevée.} Deux croquis
l'accompagnent : une sphère portant des vecteurs tangents, et un faisceau de
courbes issues d'un point, avec les étiquettes $R_i$,
$\tilde R_i = (e_j - e_k) \wedge \omega_{jk}$, $V(\omega)$,
$\varphi_\omega(\tilde R_k)$, $P_i$, dont le texte ne dit pas le
sens.\footnote{Les étiquettes sont en partie incertaines ; la ligne
« $\rho$, $\sigma$ opérant sur l'une des sections » et le mot « genres » sont de
lecture douteuse. Les vecteurs tangents à la sphère et les $e_j - e_k$
annoncent peut-être le fibré tangent de la page 19 et les sections de
$\mathcal{O}(1)$ des pages 22 et 29 ; ce n'est qu'une conjecture de lecture.}

\pagerange{15}{19}
\section*{Les groupes $\mathcal{T}_{1,1}$ et $S\mathcal{T}_{1,1}$ (pages 15, 17 et 19)}

\subsection*{$M_{1,1}$ comme orbifold de $U_{0,3}$ (page 15)}

Soit $\mathbb{D}'_3$ le produit fibré
\[
\mathbb{D}'_3 = \mathfrak{S}_3 \times_{\mathbb{Z}/2} \mathbb{Z}/4 ,
\]
au-dessus de la signature et de la réduction $\mathbb{Z}/4 \to \mathbb{Z}/2$ :
extension de $\mathbb{D}_3 = \mathfrak{S}_3$ par $\mathbb{Z}/2$, de noyau engendré
par $\omega = (1, 2)$. On fait agir $\mathbb{D}'_3$ sur $U_{0,3}$ à travers
$\mathbb{D}_3$, $\omega$ agissant trivialement. Alors (sur $\mathbb{C}$)
\[
M_{1,1} \simeq (U_{0,3}, \mathbb{D}'_3), \qquad
\mathcal{T}_{1,1} \simeq \pi_1(U_{0,3}, \mathbb{D}'_3; P_+) =
\Pi^{\mathbb{D}'}_{0,3} \simeq SL(2, \mathbb{Z}) ,
\]
l'isomorphisme dépendant du choix d'un point du demi-plan de Poincaré au-dessus
de $P_+$, par exemple $\zeta = e^{2i\pi/6}$.\footnote{La page note « à vérifier »,
barré, puis « OK ». La base n'est pas précisée ; l'énoncé est vrai sur
$\mathbb{C}$, où $M_{1,1}$ est le champ $[\mathfrak{H}/SL(2, \mathbb{Z})]$ et
$U_{0,3} = \mathfrak{H}/F$ pour le sous-groupe libre distingué $F = \ker(SL(2,
\mathbb{Z}) \to \mathbb{D}'_3)$ d'indice $12$. La page note ce groupe
$\Pi^{\mathbb{D}'_3}_{0,3}$, et écrit « $\mathbb{D}'_2 \to \mathbb{D}_2$ » là où le
contexte demande $\mathbb{D}'_3 \to \mathbb{D}_3$.} En termes de groupes :
$SL(2, \mathbb{Z})$ est l'image inverse, par $\Pi^{\mathbb{D}}_{0,3} \to
\mathbb{D}_3$, de l'extension $\mathbb{D}'_3$ de $\mathbb{D}_3$ par $\mathbb{Z}/2$,
elle-même image inverse de l'extension $\mathbb{Z}/4$ de $\mathbb{Z}/2$ par
$\mathbb{Z}/2$ ; les deux carrés

\begin{tikzcd}[column sep=small, row sep=small]
SL(2, \mathbb{Z}) \arrow[r] \arrow[d] & \Pi^{\mathbb{D}}_{0,3} \arrow[d] \\
\mathbb{D}'_3 \arrow[r] \arrow[d] & \mathbb{D}_3 \arrow[d, "\mathrm{sg}"] \\
\mathbb{Z}/4 \arrow[r] & \mathbb{Z}/2
\end{tikzcd}

sont cartésiens, avec $\rho \mapsto 0$, $\sigma \mapsto 1$ dans $\mathbb{Z}/2$.
Concrètement, $\rho' = (\rho, 2)$ et $\sigma' = (\sigma, 3)$ dans
$\Pi^{\mathbb{D}}_{0,3} \times_{\mathbb{Z}/2} \mathbb{Z}/4$ vérifient
$\rho'^3 = \sigma'^2 = (1, 2) = \omega$ et $\rho'^6 = \sigma'^4 = 1$, ce qui est la
présentation de $SL(2, \mathbb{Z})$.\footnote{Note marginale : « Bien vérifié
ainsi, notamment que $SL(2, \mathbb{Z}) \simeq \{\rho', \sigma' \mid \rho'^3 =
\sigma'^2, \ldots\}$ ». La phrase suivante de la page, en partie illisible,
décrit $SL(2, \mathbb{Z})$ comme extension de $\mathbb{D}'_3$ par un sous-groupe
isomorphe à $\Pi_{0,3}$, et s'arrête au bord de la feuille.}

\subsection*{$S\mathcal{T}_{1,1}$ (page 17)}

Soit $S\mathcal{T}_{1,1}$ l'image inverse de $SL(2, \mathbb{Z})$ dans le revêtement
universel $\widetilde{SL(2, \mathbb{R})}$. Alors
\[
S\mathcal{T}_{1,1} \simeq \langle \rho, \sigma \mid \rho^3 = \sigma^2 \rangle,
\qquad \omega = \rho^3 = \sigma^2 \text{ central}, \qquad
\mathcal{T}_{1,1} = S\mathcal{T}_{1,1}/\langle \omega^2 \rangle ,
\]
le noyau $\pi_1(SL(2, \mathbb{R})) \simeq \mathbb{Z}$ de
$\widetilde{SL(2, \mathbb{R})} \to SL(2, \mathbb{R})$ étant engendré par $\omega^2$.
C'est le groupe des tresses à trois brins, $B_3$, et le groupe du nœud de
trèfle.\footnote{Aucun de ces deux noms n'est sur la page. L'isomorphisme
$B_3 \simeq \widetilde{SL}(2, \mathbb{Z})$ et la présentation
$\langle \rho, \sigma \mid \rho^3 = \sigma^2 \rangle$ sont classiques (Artin,
1925, pour les tresses ; Schreier, 1924, pour le trèfle).} Les abélianisés sont\footnote{La page écrit $(\mathcal{T}_{11})_{\mathrm{ab}} \simeq \mathbb{Z}$ pour
le premier, sans le $S$ que demandent les valeurs dans $\mathbb{Z}$.}
\[
(S\mathcal{T}_{1,1})_{\mathrm{ab}} \simeq \mathbb{Z}, \quad
\rho \mapsto 2,\ \sigma \mapsto 3,\ \omega \mapsto 6 ;
\qquad
(\mathcal{T}_{1,1})_{\mathrm{ab}} \simeq \mathbb{Z}/12, \quad
\rho \mapsto 2,\ \sigma \mapsto 3 ,
\]
 et le carré

\begin{tikzcd}[column sep=small, row sep=small]
S\mathcal{T}_{1,1} \arrow[r] \arrow[d] & \mathcal{T}_{1,1} \arrow[d] \\
\mathbb{Z} \arrow[r] & \mathbb{Z}/12\mathbb{Z}
\end{tikzcd}

est cartésien, les deux flèches horizontales ayant le même noyau,
$\langle \omega^2 \rangle$ en haut et $12\mathbb{Z}$ en bas. Donc
$S\mathcal{T}_{1,1}$ se déduit de $\mathcal{T}_{1,1}$ par image inverse de
l'extension canonique de $\mathbb{Z}/12$ par $\mathbb{Z}$, via l'abélianisation.

Le nom est celui des groupes de Teichmüller de la \emph{Longue Marche} : pour le
tore privé d'un disque, $\mathcal{T}_{1,1}$ est le groupe des classes
d'isotopie de difféomorphismes qui préservent l'orientation, et
$S\mathcal{T}_{1,1}$ celui des difféomorphismes qui fixent le bord point par
point ; le générateur $\omega^2$ du noyau est le twist de Dehn le long du
bord.\footnote{L'interprétation de $S\mathcal{T}$ par la surface à bord est
celle de la page 20, que la présente lecture applique ici au cas $(1,1)$ ; les
pages 17 et 19 ne la formulent pas. La page 17 note, à gauche de la suite
exacte, « $\varpi = \omega^2$ ».}

\pagerange{19}{19}
\subsection*{Le diagramme de la page 19, et une question}

La page 19 rassemble ces identifications en un seul diagramme. Avec
$\tilde\rho = (\rho, 2)$, $\tilde\sigma = (\sigma, 3)$ :
\[
S\mathcal{T}_{1,1} \simeq \Pi^{\mathbb{D}}_{0,3} \times_{\mathbb{Z}/6} \mathbb{Z},
\qquad
\mathcal{T}_{1,1} \simeq \Pi^{\mathbb{D}}_{0,3} \times_{\mathbb{Z}/6} \mathbb{Z}/12
\simeq \Pi^{\mathbb{D}}_{0,3} \times_{\mathbb{Z}/2} \mathbb{Z}/4 ,
\]
où $\Pi^{\mathbb{D}}_{0,3} \to \mathbb{Z}/6$ est l'abélianisation $\rho \mapsto 2$,
$\sigma \mapsto 3$ ; tous les carrés du diagramme suivant sont cartésiens :

\begin{tikzcd}[column sep=small, row sep=small, nodes={font=\scriptsize}]
S\mathcal{T}_{1,1} \arrow[r] \arrow[d] & \mathcal{T}_{1,1} \arrow[r] \arrow[d] & \Pi^{\mathbb{D}}_{0,3} \arrow[d] \\
\mathbb{Z} \arrow[r] & \mathbb{Z}/12\mathbb{Z} \arrow[r] \arrow[d] & \mathbb{Z}/6\mathbb{Z} \arrow[d] \\
 & \mathbb{Z}/4\mathbb{Z} \arrow[r] & \mathbb{Z}/2\mathbb{Z}
\end{tikzcd}

et $\mathcal{T}_{1,1} \to \mathbb{D}'_3 = \mathbb{D}_3 \times_{\mathbb{Z}/2}
\mathbb{Z}/4$ est surjectif.\footnote{La page écrit le sommet de droite
$\Pi^{\mathbb{D}'}_{03}$ (lecture incertaine de l'exposant) ; la structure du
diagramme — un groupe d'abélianisé $\mathbb{Z}/6$, de quotient $\mathbb{D}_3$ —
demande $\Pi^{\mathbb{D}}_{0,3}$, $\Pi^{\mathbb{D}'}_{0,3}$ étant
$\mathcal{T}_{1,1}$ lui-même. Elle écrit aussi, à droite, « $\mathcal{T}_{11}$
identifié à $\Pi^{D}_{03} \times_{\mathbb{Z}/12\mathbb{Z}} \mathbb{Z}/4\mathbb{Z}$ »
(le $4$ incertain) : $\Pi^{\mathbb{D}}_{0,3}$ n'a pas de surjection sur
$\mathbb{Z}/12$, et les formes justes sont celles qu'on vient d'écrire. Les
positions de deux des marques « cart » sont incertaines.}

Puis la question : « A-t-on bien $S\mathcal{T}_{11} \simeq S\Pi^{D}_{03}$ ?? », à
laquelle la page 33 répondra oui. Et une conjecture :
« il semblerait plutôt $PM_{11} \simeq (T^{*}U_{03}, \mathfrak{S}'_3)$ (fibré
cotangent de $U_{03}$ !!), donc $\mathcal{T}_{11} \simeq \pi_1(T^{*}U_{03},
\mathfrak{S}'_3)$ ». Le symbole $PM_{1,1}$ n'est défini nulle part ; la page 26
en donne des variantes de niveau $2$, et les pages 34 et 35 le remplacent par le
$\mathbb{G}_m$-torseur $PU_{0,3} = \mathcal{O}(1)^{*}$.\footnote{Si $PM_{1,1}$
désigne, comme le suggère la page 26, le $\mathbb{G}_m$-torseur des couples
(courbe elliptique, différentielle invariante non nulle), son groupe fondamental
est $S\mathcal{T}_{1,1}$ et non $\mathcal{T}_{1,1}$ : cet espace est le
complémentaire de la cubique cuspidale $g_2^3 = 27g_3^2$ dans $\mathbb{C}^2$, de
groupe fondamental $B_3$. Le lien entre le carré de la différentielle invariante
et le fibré cotangent de l'espace des modules est ce qu'on appelle
l'isomorphisme de Kodaira--Spencer ; la page ne le nomme pas, et la lecture ne
prétend pas que ce soit la raison du « !! ». « tangent » est biffé et remplacé
par un mot lu « cotangent » d'après le $T^{*}$.}

\pagerange{20}{20}
\section*{Surfaces à bord : $S\mathcal{T}$, $\mathcal{T}$ et les tours de bord (page 20)}

La page 20 pose le cadre général dont $\mathcal{T}_{1,1}$ et $S\mathcal{T}_{1,1}$
sont le premier cas. Soit $X$ une surface compacte orientée à $n$ composantes de
bord, $\pi$ son groupe fondamental et $\ell_1, \dots, \ell_n$ des lacets autour
des composantes du bord. Un automorphisme de la surface induit
$u \in \mathrm{Aut}\,\pi$, une permutation $\sigma \in \mathfrak{S}_n$ et des
éléments $g_i \in \pi$ tels que
\[
u(\ell_i) = g_i\,\ell_{\sigma(i)}\,g_i^{-1} ,
\]
et la page forme le groupe $\widetilde{E}$ de ces données, qui contient $\pi$
(par les automorphismes intérieurs) et $\mathbb{Z}^n$ (par les changements
$g_i \mapsto g_i\ell_{\sigma(i)}^{k_i}$), avec $\widetilde{E} \to \mathfrak{S}_n$ ;
on pose $E = \widetilde{E}/\pi$.\footnote{Cette lecture des données du cadre de la
page, et du plongement $\mathbb{Z}^n \hookrightarrow \widetilde{E}$ (tracé avec un
crochet à chaque bout), est de la présente lecture ; la page ne les commente
pas. Elle écrit aussi une flèche $Z\pi \to \mathbb{Z}^n$.} La fibration
\[
S\mathrm{Aut}^{+}(X) \to \mathrm{Aut}^{+}(X) \to \mathrm{Aut}^{+}(\partial X)
\simeq \mathfrak{S}_n \ltimes \mathbb{T}^n ,
\]
où $S\mathrm{Aut}^{+}(X)$ est le groupe des difféomorphismes qui fixent le bord
point par point et où l'équivalence de droite est une équivalence d'homotopie
($\mathrm{Diff}^{+}(S^1) \simeq SO(2)$), donne, avec
$S\mathcal{T} = \pi_0 S\mathrm{Aut}^{+}(X)$ et $\mathcal{T} = \pi_0
\mathrm{Aut}^{+}(X)$, la suite exacte
\[
1 \to \pi_1\mathrm{Aut}^{+}(X) \to \mathbb{Z}^n \to S\mathcal{T} \to \mathcal{T}
\to \mathfrak{S}_n \to 1 ,
\]
où $\mathbb{Z}^n = \pi_1(\mathbb{T}^n)$ s'envoie dans $S\mathcal{T}$ par les
twists de Dehn le long des composantes du bord, et $\pi_1\mathrm{Aut}^{+}(X)$ vaut
$0$ ou $\mathbb{Z}$.\footnote{La suite commence par $1$ parce que
$\pi_1 S\mathrm{Aut}^{+}(X) = 0$, les composantes neutres de ces groupes étant
contractiles dès que la caractéristique d'Euler est négative (Earle--Eells 1969,
Earle--Schatz 1970) ; la page ne le dit pas. Le cas $\mathbb{Z}$ est celui du
disque et de l'anneau. La page note $\mathfrak{S}_n \cdot \mathbb{T}^n$ le produit
en couronne, et écrit encore un groupe $\mathrm{Aut}^{+}_{!}(\widetilde{X})$ que
rien ne définit.} Elle écrit encore $1 \to S\mathcal{T} \to E \to \mathfrak{S}_n
\to 1$.

La fin de la page abstrait la situation. Soit $1 \to F \to A \to G \to 1$ une
extension de groupes topologiques, $F$ discret. La suite d'homotopie de la
fibration $F \to A \to G$ s'écrit
\[
1 \to \pi_1(A) \to \pi_1(G) \to F \to \pi_0(A) \to \pi_0(G) \to 1 ,
\]
l'image de $\pi_1(G)$ dans $F$ est centrale — c'est $F \cap A^{\circ}$, et la
composante neutre, connexe, agit trivialement par conjugaison sur le sous-groupe
distingué discret $F$ —, et $\pi_0(G)$ opère extérieurement sur $F$, donc sur
son centre $Z(F)$. En posant $E = \pi_0(\widetilde{A})$ pour une extension
$\widetilde{A}$ convenable, on obtient $1 \to F \to E \to \pi_0 \to
1$.\footnote{« central » est de lecture incertaine ; l'argument qui le justifie est
de la présente lecture. Le diagramme de la page met en regard deux extensions
$\widetilde{F} \to \widetilde{A} \to \widetilde{G}$ et $F \to A \to G$, $F$
marqué « discret » et $G = \mathfrak{S}_n \cdot \mathbb{T}^n$ ; la définition de
$\widetilde{A}$ n'est pas écrite.}

\pagerange{16}{16}
\section*{Cartes : sommets, arêtes, faces, drapeaux (pages 16, 29 et 31)}

\subsection*{Polyèdres réguliers et drapeaux (page 16)}

Pour un polyèdre $X_\bullet$ dont le groupe $G_0$ des automorphismes agit
simplement transitivement sur un ensemble $\vec{S}(X_\bullet)$, les ensembles de
sommets, d'arêtes et de faces sont des espaces homogènes
\[
S(X_\bullet) \simeq G_0/H_S, \qquad A(X_\bullet) \simeq G_0/H_A, \qquad
F(X_\bullet) \simeq G_0/H_F ,
\]
et les stabilisateurs d'un sommet, d'une arête et d'une face incidents se coupent
deux à deux trivialement, ce qui équivaut aux injections\footnote{La page écrit deux fois la deuxième injection ; la troisième, que
demande $H_S \cap H_F = e$, est rétablie. Elle ne définit pas $\vec{S}$ ; on lit
un ensemble principal homogène sous $G_0$, par exemple celui des sommets munis
d'une arête issue d'eux (les « brins »), lecture qui est la nôtre.}
\[
G_0 \hookrightarrow G_0/H_S \times G_0/H_A, \qquad
G_0 \hookrightarrow G_0/H_A \times G_0/H_F, \qquad
G_0 \hookrightarrow G_0/H_S \times G_0/H_F .
\]


Dans le cas non orienté, pour la catégorie des polyèdres d'un type donné $\tau$,
on a les foncteurs $S$, $A$, $F$, les relations d'incidence
$A_1 \subset S \times A$, $A_2 \subset A \times F$, $A_3 \subset S \times F$, et
l'ensemble des \emph{drapeaux}
\[
R = A_1 \times_A A_2 = \{ (s, a, f) \mid s < a < f \}, \qquad
A_1 = \mathrm{Im}(R \to S \times A),\ \ldots,\ A_3 = \mathrm{Im}(R \to S \times F) .
\]
Avec $\Omega$ l'ensemble des orientations, la page écrit $A_{1\,\Omega} \simeq
A_{2\,\Omega} \simeq A_{3\,\Omega}$ et $\mathrm{Pol}_\tau^{\Omega} \simeq
\mathrm{Pol}_{\tau/\Omega}$, sans les définir.\footnote{La parenthèse qui suit
$A_{3\,\Omega}$ reste ouverte ; l'exposant $\Omega$ de $\mathrm{Pol}_\tau$ est
écrit sur une rature.} C'est la description des cartes par l'action d'un groupe
sur les drapeaux, que l'\emph{Esquisse d'un programme} (1984) appellera le groupe
cartographique.\footnote{Le mot n'est pas dans le dossier.}

\pagerange{29}{29}
\subsection*{Le changement d'orientation (page 29, moitié inférieure)}

Sur l'ensemble $\bar A$ des arêtes d'une carte combinatoire orientée, une
orientation définit deux opérations, $\vec\rho_s$ (autour des sommets) et
$\vec\rho_f$ (autour des faces), liées par
\[
\vec\rho_s^{\,-1}\,\vec\rho_f = \sigma ,
\]
$\sigma$ l'involution des arêtes. La page calcule ce qu'elles deviennent pour
l'orientation opposée : si $\bar a$ provient de $a$ et donc aussi de
$\sigma_n a$, on a $\rho_s(\sigma_n a) = \sigma_n(\rho_s^{-1}a)$, d'où
\[
\vec\rho\,'_s = \vec\rho_s^{\,-1}, \qquad
\vec\rho\,'_f = \vec\rho\,'_s\,\sigma = \vec\rho_s^{\,-1}\sigma
= \sigma\,\vec\rho_f^{\,-1}\sigma = \mathrm{int}(\sigma)\bigl(\vec\rho_f^{\,-1}\bigr)
.
\]
Renverser l'orientation inverse les deux permutations, celle des faces à
conjugaison par $\sigma$ près.\footnote{La page écrit
« $\vec\rho\,'_s = \vec\rho_f^{\,-1}$ », ce qui contredit la ligne précédente ;
la suite utilise $\vec\rho_s^{\,-1}$, qui est juste. La moitié inférieure de la
page est écrite tête-bêche ; « on a vu » renvoie à des pages qui ne sont pas dans
ce dossier, et le nom de l'objet (« complexe » biffé, puis un mot illisible,
« combinatoire ») n'est pas sûr.}

\pagerange{31}{31}
\subsection*{Orbites de $\rho$ et triples $(A, \sigma, \rho)$ (page 31)}

La page 31 construit les orbites d'une opération $\rho$ sur un ensemble $A$ d'arêtes
orientées pas à pas, et distingue deux cas : ou bien la construction s'arrête —
on finit par sortir de la partie $A'$ où $\rho$ est défini, et l'on ferme le
cycle —, ou bien elle se prolonge indéfiniment, et l'on ne définit pas $\rho$.
Supposant désormais $\rho$ défini partout, $\rho$ est une permutation de $A$ ;
$\bar S = A/\rho$ est l'ensemble de ses orbites, $\hat S \subset \bar S$ celui des
orbites finies, et l'on compactifie partiellement la surface $X$ en $\hat X$ en
lui ajoutant les points de $\hat S - S$. La conclusion est que les quadruples
$(A, \sigma, A', \rho)$ satisfaisant aux axiomes, avec $A' = \hat A'$,
s'identifient aux triples
\[
(A, \sigma, \rho), \qquad A \text{ fini},\ \sigma \text{ involution sans point
fixe},\ \rho \text{ permutation quelconque} ,
\]
$A'$ se retrouvant comme une partie de $A$ telle que $A' \cup \sigma A' =
A$.\footnote{La page est une rédaction rapide dont beaucoup de mots de liaison
sont illisibles, et plusieurs phrases ne se laissent pas reconstruire ; les
axiomes des quadruples sont dans des pages absentes du dossier. « fini » est de
lecture incertaine.} C'est la description moderne d'une carte orientée par ses
brins : $\sigma$ recolle les deux moitiés de chaque arête, $\rho$ tourne autour
des sommets, et les faces sont les orbites de $\rho\sigma$. Le groupe qui agit est
$\langle \rho, \sigma \mid \sigma^2 = 1 \rangle$, et les cartes dont tous les
sommets sont de valence $3$ sont les ensembles sur lesquels agit
$\langle \rho, \sigma \mid \rho^3 = \sigma^2 = 1 \rangle = \Pi^{\mathbb{D}}_{0,3}$ :
c'est ce qui relie ces pages au reste du dossier.\footnote{Ce lien est de la
présente lecture ; il est exposé dans l'\emph{Esquisse d'un programme}, § 3, et
non dans ce dossier.}

\pagerange{22}{24}
\section*{Le lemme : $H^1(U_{0,3}, \mathfrak{S}_3; \mathbb{G}_m) \simeq \mathbb{Z}/6$ (pages 22, 24 et 29)}

\subsection*{L'énoncé (page 22)}

Soit $J$ un ensemble à trois éléments, $S$ un schéma, et
\[
0 \to V_J \to \mathcal{O}_S^{J} \to \mathcal{O}_S \to 0 ,
\]
$V_J$ le module libre de rang $2$ des combinaisons de somme nulle des $e_i$,
$i \in J$ ; soit $\Sigma_J = \mathbb{P}(V_J)$, droite projective sur $S$ munie de
l'action de $\mathfrak{S}_J \simeq \mathfrak{S}_3$, avec
$\Gamma(\Sigma_J, \mathcal{O}(1)) = V_J$, engendré par les $e_i - e_j$
(page 29). Les trois sections $e_i - e_j$ s'annulent en trois points ; leur
complémentaire $\Sigma_J^{*}$ est une copie de $U_{0,3}$ sur laquelle
$\mathfrak{S}_3$ agit par les substitutions anharmoniques. On note
$H^1(U_{0,3}, \mathfrak{S}_3; \mathbb{G}_m)$ le groupe des classes de
$\mathbb{G}_m$-torseurs $\mathfrak{S}_3$-équivariants sur $U_{0,3}$ — le groupe de
Picard du champ $(U_{0,3}, \mathfrak{S}_3)$.

\textbf{Lemme.} \emph{Soit $S$ un schéma régulier intègre, avec
$\mathrm{Pic}(S) = 0$ et de caractéristique générique $\neq 2$. Alors
\[
H^1(U_{0,3}, \mathfrak{S}_3; \mathbb{G}_m) \simeq \mathbb{Z}/6\mathbb{Z},
\]
engendré par la classe de $\mathcal{O}(1)$, c'est-à-dire du
$\mathbb{G}_m$-torseur $P\Sigma_J = \mathbb{V}(\mathcal{O}(1))^{*}$, restreint à
$\Sigma_J^{*}$. En caractéristique générique $2$, on trouve
$\mathbb{Z}/3\mathbb{Z}$.}

\emph{Majoration} (note marginale). La suite spectrale de Hochschild--Serre donne
\[
0 \to H^1\bigl(\mathfrak{S}_3, H^0(U_{0,3}, \mathbb{G}_m)\bigr) \to
H^1(U_{0,3}, \mathfrak{S}_3; \mathbb{G}_m) \to
H^0\bigl(\mathfrak{S}_3, H^1(U_{0,3}, \mathbb{G}_m)\bigr) = 0 ,
\]
le dernier terme étant nul parce que $\mathrm{Pic}(U_{0,3}) = 0$ sous les
hypothèses faites sur $S$ ; et
\[
0 \to H^0(S, \mathbb{G}_m) \to H^0(U_{0,3}, \mathbb{G}_m) \to V(J) \to 0 ,
\]
où $V(J) \simeq \ker(\mathbb{Z}^J \xrightarrow{\Sigma} \mathbb{Z})$ est le groupe
des diviseurs des unités, portés par les trois points retirés. Or
$H^1(\mathfrak{S}_3, \mathcal{O}(S)^{*}) = \mathrm{Hom}(\mathfrak{S}_3,
\mathcal{O}(S)^{*}) = \mu_2(S)$, d'ordre $2$ si la caractéristique générique est
$\neq 2$ et $1$ sinon, et $H^1(\mathfrak{S}_3, V(J)) = \mathbb{Z}/3$ (suite exacte
de $0 \to V(J) \to \mathbb{Z}^J \to \mathbb{Z} \to 0$, avec
$H^1(\mathfrak{S}_3, \mathbb{Z}^J) = H^1(\mathfrak{S}_2, \mathbb{Z}) = 0$). Le groupe
cherché est donc d'ordre divisant $6$ (resp. $3$), et il suffit de montrer que la
classe de $\mathcal{O}(1)$ est d'ordre exactement $6$ (resp. $3$).\footnote{La
note marginale écrit les deux suites exactes, et les calculs de
$H^1(\mathfrak{S}_3, \cdot)$ ne sont pas sur la page ; ils sont ajoutés. Le
terme médian de la première suite est coupé par le bord de la feuille.}

\emph{L'ordre divise $6$.} La section
\[
\varphi = \prod_{\substack{i, j \in J \\ i \neq j}} (e_i - e_j)
= -\prod_{\{j, k\} \subset J} (e_j - e_k)^2 = -\Delta
\ \in \Gamma(\Sigma_J, \mathcal{O}(6)) = \mathrm{Sym}^6 V_J
\]
est invariante par $\mathfrak{S}_3$ et ne s'annule pas sur $\Sigma_J^{*}$ : elle
trivialise $\mathcal{O}(6)$ de façon équivariante. Le groupe structural de
$P\Sigma_J^{*}$ se réduit ainsi à $\mu_6$ : c'est le $\mu_6$-torseur
$\{\lambda \mid \lambda^6 = \varphi\}$, que la page 29 écrit
$\lambda^6 = \prod_{i,j}\xi_{ij}(x)$.\footnote{La page écrit
$\varphi = \prod_{i \in J} (e_j - e_k)^2$, oubliant le signe $(-1)^3$ ; un $\pm$
biffé devant le produit et le « NB Signe erroné » qui suit montrent qu'elle l'a
vu. Le signe est sans effet : $-\Delta$ est invariant comme $\Delta$. La ligne
« $\mathcal{O}(6) \simeq V_J^{\otimes 3}(\omega)$ » et la phrase qui l'entoure,
très surchargées, ne se lisent qu'en partie.}

\emph{Les invariants.} Lorsque $6$ est inversible sur $S$, l'anneau des
invariants est un anneau de polynômes,
\[
\mathrm{Sym}^{*}(V_J)^{\mathfrak{S}_3} = \mathcal{O}[\gamma_2, \gamma_3], \qquad
\gamma_2 = \sum_{i \in J} (e_j - e_k)^2, \qquad
\gamma_3 = \prod_{i \in J} (e_j + e_k - 2e_i) ,
\]
et le discriminant s'y écrit $54\,\Delta = \gamma_2^{3} - 2\gamma_3^{2}$. La courbe
$\gamma_2^3 = 2\gamma_3^2$ du plan des invariants est une cubique cuspidale ; son
complémentaire est le quotient, libre, par $\mathfrak{S}_3$ du plan privé des
trois droites où s'annulent les $e_i - e_j$, c'est-à-dire de $P\Sigma_J^{*}$.\footnote{La page écrit
« $\mathrm{Sym}^{*}(V_J)^{\mathfrak{S}_3} \simeq
\mathcal{O}[\gamma_2, \gamma_3]/(2\gamma_3^2 - \gamma_2^3)$ » et
« $\Delta = \gamma_2^3$ », de lecture difficile. Telle quelle, la première
égalité est fausse : $\gamma_2$ et $\gamma_3$ sont algébriquement indépendants
(théorème de Chevalley pour le groupe de réflexions $\mathfrak{S}_3$), et
$\gamma_2^3 - 2\gamma_3^2$ n'est pas une relation mais, à un facteur près, le
discriminant. On écrit ce qui est vrai. Pour le vérifier, on pose
$x_i = e_i - \frac13\sum e$ ; alors $\gamma_2 = -6p$, $\gamma_3 = -27q$ avec
$p = \sum x_jx_k$, $q = x_1x_2x_3$, et $\Delta = -4p^3 - 27q^2$. L'hypothèse
« $6$ inversible » est ajoutée ; le lemme lui-même n'en a pas besoin.}

\pagerange{24}{24}
\subsection*{L'ordre est exactement $6$ (page 24)}

$\gamma_2$ s'annule aux deux points $P_\pm$ de $\Sigma_J^{*}$ fixés par les
éléments d'ordre $3$, et $\gamma_3$ aux trois points $Q_i$ fixés par les
transpositions. Il n'y a donc pas de section inversible invariante de
$\mathcal{O}(n)$ sur $\Sigma_J^{*}$ pour $0 < n < 6$, et la classe de
$\mathcal{O}(1)$ est d'ordre $6$ : c'est un générateur.\footnote{La page écrit
« pour $n \leqslant 6$ » ; l'argument demande $n < 6$, puisque $\varphi$ est une
telle section pour $n = 6$. L'identification des $P_i$ et des $Q_i$ aux points à
stabilisateur d'ordre $3$ et $2$ est de la présente lecture ; la page ne les
définit pas, et les $P_\pm$ sont vraisemblablement le $P_+$ des pages 7 et 15.}

Le raisonnement se dit mieux ainsi : en un point fixe d'un sous-groupe $H$, $H$
agit sur la fibre de $\mathcal{O}(n)$ par un caractère ; une section invariante
qui ne s'annule pas en ce point exige que ce caractère soit trivial. En $P_\pm$ le
stabilisateur, d'ordre $3$, agit sur la fibre de $\mathcal{O}(1)$ par une racine
primitive cubique de l'unité, en $Q_i$ le stabilisateur, d'ordre $2$, par $-1$ :
il faut $3 \mid n$ et $2 \mid n$. Un argument qui vaut aussi en caractéristique
$3$, où $P_+ = P_-$ : une section inversible de $\mathcal{O}(n)$ sur
$\Sigma_J^{*}$ est, à une unité de $S$ près, de la forme
$\prod (e_j - e_k)^{a_i}$ avec $\sum a_i = n$ ; l'invariance par les
permutations circulaires impose $a_1 = a_2 = a_3$, et celle par les
transpositions, qui changent le signe de $\prod_{j<k}(e_j - e_k)$, que cet
exposant commun soit pair ; d'où $6 \mid n$. En caractéristique $2$, le signe
disparaît, $\prod_{j<k}(e_j - e_k)$ est invariante, et l'ordre est $3$.\footnote{Ces
deux formes de l'argument sont de la présente lecture.}

Le résultat est l'analogue, pour le champ $(U_{0,3}, \mathfrak{S}_3)$, du calcul
$\mathrm{Pic}(M_{1,1}) \simeq \mathbb{Z}/12$ de Mumford (1965) : le passage de
$\mathbb{D}_3$ à $\mathbb{D}'_3$ double l'ordre.\footnote{Le rapprochement est de
la présente lecture ; le dossier ne cite pas Mumford.}

\pagerange{29}{29}
\subsection*{Les sections $e_i - e_j$ (page 29, moitié supérieure)}

La page 29 revient sur ces sections : $V_J = \Gamma(X, \mathcal{O}(1))$ est
engendré par les $e_i - e_j$, que les permutations échangent entre elles au signe
près ; une fonction $\lambda$ avec $\lambda^6 = \prod_{i \neq j}\xi_{ij}(x)$, où
$\xi_{ij}$ exprime $e_i - e_j$ dans une trivialisation locale ; la suite
$\mathbb{Z} \to \widetilde{\Pi}_{0,3} \to \Pi_{0,3} \to 1$ ; et
« $V(\omega) \overset{!!}{\simeq} V$ ».\footnote{La prose du haut de la page est
presque entièrement illisible ; la lecture ne retient que les formules sûres.
Quelques nombres ($\pm\frac12$, $\pm 1$, $(\frac12)^4$, $\sqrt[3]{1/2}$) sont
dispersés sans relation écrite, et une colonne en marge porte
« $\tilde R_K - \tilde R_j$ ? ».}

\pagerange{26}{26}
\section*{Les champs au-dessus de $M_{1,1}$ (page 26)}

La page 26 est faite de diagrammes sans texte. On les lit dans le langage des
champs, avec $\omega$ l'automorphisme $-1$ d'une courbe elliptique :
$(U_{0,3}, \{1, \omega\})$, le produit de $U_{0,3}$ par la gerbe triviale
$B\mathbb{Z}/2$, est le champ $M_{1,1}(2)$ des courbes elliptiques munies d'une
structure de niveau $2$ ; $(U_{0,3}, \mathbb{D}'_3) \simeq M_{1,1}$ (page 15) ; et
$(U_{0,3}, \mathbb{D}_3)$, où l'on a oublié $\omega$, est ce qu'on appelle
aujourd'hui la rigidification de $M_{1,1}$ le long de $\mu_2$, que la page note
$M'_{1,1}$. Le premier diagramme en donne les flèches, chacune marquée de son
groupe ou de sa nature :

\begin{tikzcd}[column sep=small, row sep=small, nodes={font=\scriptsize}]
U_{0,3} \arrow[r, "{\{1, \omega\}}"] \arrow[dr, "\mathbb{D}'_3"'] & (U_{0,3}, \{1, \omega\}) \arrow[d, "\mathbb{D}_3"'] \arrow[dr, "\mathrm{gerbe}"] & \\
 & (U_{0,3}, \mathbb{D}'_3) \arrow[d, "\mathrm{gerbe}"'] & U_{0,3} \arrow[dl, "\mathbb{D}_3"] \\
 & (U_{0,3}, \mathbb{D}_3) &
\end{tikzcd}

— les flèches marquées d'un groupe sont des revêtements galoisiens de ce groupe,
les flèches « gerbe » des gerbes de lien $\mathbb{Z}/2$. Le second diagramme est
le même, calqué sur le demi-plan de Poincaré $\mathcal{D}$ : $(\mathcal{D},
\Pi_{0,3}) \simeq \mathcal{D}/\Pi_{0,3}$, $(\mathcal{D}, \Pi_{0,3} \times \{1,
\omega\})$, $(\mathcal{D}, SL(2, \mathbb{Z}))$ et $(\mathcal{D}, PSL(2,
\mathbb{Z}))$, les flèches étant marquées « gal » et
« 2-gerbe ».\footnote{Redessiné sur une grille : sur la page, la flèche marquée
$\mathbb{D}_3$ descend en oblique et la flèche « gerbe » verticalement. La page
écrit $SL(2, \mathbb{Z})'$ pour $PSL(2, \mathbb{Z})$ ; les étiquettes des
flèches « gal » sont surchargées ($\Pi^{\mathbb{D}'}_{0,3}$ ou $\mathbb{D}'_3$,
$\mathbb{D}_3$) ; la lecture de $\mathcal{D}$ comme demi-plan de Poincaré est
incertaine.} Un troisième diagramme relie $M_{1,1}(\tilde 2)$, $M_{1,1}(2)$,
$M_{1,1}(2)'$, $M_{1,1}$ et $M'_{1,1}$, et le carré cartésien
\[
\begin{array}{ccc}
PM_{1,1}(\tilde 2) \simeq PU_{0,3} & \longrightarrow & PM_{1,1}(2) \simeq TU_{0,3} \\
\downarrow & & \downarrow \\
U_{0,3} \simeq M_{1,1}(\tilde 2) & \longrightarrow & (U_{0,3}, \{1, \omega\}) \simeq M_{1,1}(2)
\end{array}
\]
met pour la première fois en scène le $\mathbb{G}_m$-torseur $PU_{0,3}$ des
pages 34 et 35.\footnote{Ni $M_{1,1}(\tilde 2)$, ni $PM_{1,1}$, ni $TU_{0,3}$ ne
sont définis par la page, et la lecture ne les définit pas à sa place. La flèche
du haut est tracée double ; une formule biffée, « $x \mapsto x^2 \wedge
\omega_0$ », suggère qu'elle est l'élévation au carré. Une ligne isolée porte
« $PU(TU_{03}, \mathbb{D}_3; \mu_2)$ », le $\mu_2$ incertain.}

\pagerange{28}{30}
\section*{Le diagramme des extensions (pages 28 et 30)}

Les pages 28 et 30 dessinent en perspective, puis à plat, un diagramme qui
superpose les extensions déjà rencontrées. Ses sommets finis sont quatre groupes,
tous produits fibrés au-dessus de $\mathbb{Z}/2$ (signature d'un côté, réduction
de l'autre) :
\[
\begin{array}{lll}
\mathbb{D}_3 = \mathfrak{S}_3 & \text{ordre } 6, & \text{quotient de }
\Pi^{\mathbb{D}}_{0,3} \text{ par } \Pi_{0,3} ; \\[2pt]
\mathbb{D}'_3 = \mathbb{D}_3 \times_{\mathbb{Z}/2} \mathbb{Z}/4 & \text{ordre } 12,
& \text{quotient de } SL(2, \mathbb{Z}) \text{ par } \Pi_{0,3} ; \\[2pt]
\mathbb{D}''_3 = \mathbb{D}_3 \times_{\mathbb{Z}/2} \mathbb{Z}/6 & \text{ordre } 18,
& \text{quotient de } \Pi^{\mathbb{D}}_{0,3} \text{ par } \Pi'_{0,3} ; \\[2pt]
\mathbb{D}'''_3 = \mathbb{D}_3 \times_{\mathbb{Z}/2} \mathbb{Z}/12
\simeq \mathbb{D}'_3 \times_{\mathbb{Z}/2} \mathbb{Z}/6 & \text{ordre } 36, &
\text{quotient de } SL(2, \mathbb{Z}) \text{ par } \Pi'_{0,3} .
\end{array}
\]
Le dernier isomorphisme vient de $\mathbb{Z}/12 \simeq \mathbb{Z}/4
\times_{\mathbb{Z}/2} \mathbb{Z}/6$. Les flèches $\mathbb{D}'''_3 \to
\mathbb{D}''_3$ et $\mathbb{D}'_3 \to \mathbb{D}_3$ sont de degré $2$, les flèches
$\mathbb{D}'''_3 \to \mathbb{D}'_3$ et $\mathbb{D}''_3 \to \mathbb{D}_3$ de degré
$3$, et la face formée par ces quatre groupes est cartésienne, au-dessus de la
face $\mathbb{Z}/12$, $\mathbb{Z}/6$, $\mathbb{Z}/4$, $\mathbb{Z}/2$.\footnote{La
page 28 écrit les produits fibrés au-dessus de $\mathbb{Z}/12\mathbb{Z}$
(« $\mathbb{D}_3 \times_{\mathbb{Z}/12} \mathbb{Z}/12 \simeq \mathbb{D}'_3
\times_{\mathbb{Z}/12} \mathbb{Z}/6$ », « $\mathbb{D}_3 \times_{\mathbb{Z}/12}
\mathbb{Z}/6 = E$ ») ; $\mathbb{D}_3$ n'ayant pas de surjection sur
$\mathbb{Z}/12$, seuls les produits au-dessus de $\mathbb{Z}/2$ ont un sens, et
ce sont eux qui ont les ordres $18$ et $36$ que la page inscrit. Le groupe
$\mathbb{D}'''_3$ est le $\mathcal{E}$ de la page 9 (voir plus haut). Les chiffres
$2$ et $3$ portés par les flèches sont de sa main.} Au-dessus de ces quotients
finis, la page empile les extensions infinies : dans une colonne
$\widetilde{\Pi}{}^{\circ}_{0,3} \subset \widetilde{\Gamma}(2) \subset
\widetilde{SL}(2, \mathbb{Z})$, extensions par $\mathbb{Z}$ de $\Pi_{0,3} \subset
\Gamma(2) \subset SL(2, \mathbb{Z})$, elles-mêmes au-dessus de
$\Pi_{0,3} \subset \Pi_{0,3} \subset \Pi^{\mathbb{D}}_{0,3}$ ; dans une ligne,
une extension $\widetilde{\mathbb{D}}_3$ de $\mathbb{D}'''_3$ par $\mathbb{Z}$,
quotient de $\widetilde{SL}(2, \mathbb{Z})$ ; et, en bas, l'abélianisation
$\widetilde{SL}(2, \mathbb{Z}) \to \mathbb{Z} \to \mathbb{Z}/12 \to
\mathbb{Z}/6$, le noyau de $\mathbb{Z} \to \mathbb{Z}$ étant $12\mathbb{Z}$. On
peut lire $\widetilde{\mathbb{D}}_3$ comme le produit fibré $\mathbb{D}'''_3
\times_{\mathbb{Z}/12} \mathbb{Z}$, quotient de $\widetilde{SL}(2, \mathbb{Z})$
par $(g \bmod \Pi'_{0,3},\ g_{\mathrm{ab}})$.\footnote{Le dessin de la page 28 est
une perspective de deux cubes, redessinée sur une grille dans la transcription, et
plusieurs flèches y ont une place incertaine ; on n'en retient que la structure.
Les noyaux y sont notés $\mathbb{Z}(\frac{\omega}{2})$ et $\mathbb{Z}(\omega)$,
reliés par une flèche double entourée ; la lecture ne les réécrit pas, faute de
savoir à quel $\omega$ ils se réfèrent. La description de
$\widetilde{\mathbb{D}}_3$ comme produit fibré est de la présente lecture.}

Au bas de la page 28, l'extension $\widetilde{\Pi}{}^{\mathbb{D}}_{0,3}(4)$ est
la somme amalgamée de $\widetilde{\Pi}{}^{\mathbb{D}}_{0,3}$ le long de
$\mathbb{Z} \xrightarrow{4} \mathbb{Z}$ :
\[
1 \to \mathbb{Z} \to \widetilde{\Pi}{}^{\mathbb{D}}_{0,3}(4) \to
\Pi^{\mathbb{D}}_{0,3} \to 1, \qquad\text{et aussi}\qquad
1 \to \mathbb{Z}\omega \to \widetilde{\Pi}{}^{\mathbb{D}}_{0,3}(4) \to
\Pi^{\mathbb{D}}_{0,3} \times \mathbb{Z}/4 \to 1 ,
\]
la seconde parce que l'image de l'ancien noyau est d'indice $4$ dans le nouveau, et
que l'extension quotient de $\Pi^{\mathbb{D}}_{0,3}$ par $\mathbb{Z}/4$ est
scindée. La page demande alors ce qu'est l'image inverse de $SL(2, \mathbb{Z})
\subset \Pi^{\mathbb{D}}_{0,3} \times \mathbb{Z}/4$, et écrit un point
d'interrogation : la réponse est aux pages 34 et 35.

\pagerange{30}{30}
\subsection*{La même chose, pour un $\Pi$ quelconque (page 30)}

La page 30 refait ce diagramme avec des lettres : $\Pi$, son extension
$\widetilde{\Pi}$ par $\mathbb{Z}$, un sous-groupe distingué $\Pi_0$ de quotient
$\mathbb{D}_3$, et $\Pi' \subset \Pi \times \mathbb{Z}/4$ l'image inverse de
$\mathbb{D}'_3 \subset \mathbb{D}_3 \times \mathbb{Z}/4$ :
\[
\begin{array}{ccccccccc}
1 & \to & \Pi_0 & \to & \Pi \times \mathbb{Z}/4 & \to & \mathbb{D}_3 \times
\mathbb{Z}/4 & \to & 1 \\
 & & \parallel & & \cup & & \cup & & \\
1 & \to & \Pi_0 & \to & \Pi' & \to & \mathbb{D}'_3 & \to & 1
\end{array}
\]
et l'extension $\widetilde{\Pi}(4) \to \Pi \times \mathbb{Z}/4$ par
$\mathbb{Z}$.\footnote{La page écrit « $\Pi' \subset \Pi$ », là où le diagramme
qu'elle trace demande $\Pi' \subset \Pi \times \mathbb{Z}/4$ ; la flèche de
droite, de $\mathbb{D}'_3$ vers $\mathbb{D}_3 \times \mathbb{Z}/4$, est tracée avec
une tête à chaque bout, et on la lit comme l'inclusion. Les $\mathbb{D}$ y sont
écrits comme des $D$ ordinaires.} C'est, avec $\Pi = \Pi^{\mathbb{D}}_{0,3}$ et
$\Pi_0 = \Pi_{0,3}$, la description $\Pi' = \mathcal{T}_{1,1} = SL(2, \mathbb{Z})$
de la page 15.

\pagerange{32}{34}
\section*{Classes d'extensions, et le caractère du revêtement (pages 32 et 34)}

\subsection*{Un calcul dans $\mathrm{Ext}^1$ (page 32)}

Pour $n \geqslant 1$, soit $\eta_n \in \mathrm{Ext}^1(\mathbb{Z}/n, \mathbb{Z})
\simeq \mathbb{Z}/n$ la classe de l'extension canonique
$0 \to \mathbb{Z} \xrightarrow{n} \mathbb{Z} \to \mathbb{Z}/n \to 0$. Soient $p$,
$q$ deux entiers, $\alpha \in (\mathbb{Z}/p)^{*}$, et $\widetilde{\Pi}_0(q)$
l'extension de $\mathbb{Z}/p \times \mathbb{Z}/q$ par $\mathbb{Z}$ de classe\footnote{La page construit la première composante comme l'extension
$0 \to \mathbb{Z} \xrightarrow{p} \mathbb{Z} \xrightarrow{\alpha\varphi_p}
\mathbb{Z}/p \to 0$, $\varphi_p$ la projection canonique, et lui donne la classe
$\alpha\eta_p$. Avec cette surjection la classe est en fait $\alpha^{-1}\eta_p$
(un relèvement de $1$ est un $e$ avec $\alpha e \equiv 1$, et $pe$ vaut
$\alpha^{-1}$ fois l'image du générateur) ; on définit donc l'extension par sa
classe. La différence est sans effet dans le seul cas utilisé, $\alpha = -1$.
L'argument de $\mathrm{pr}_2^{*}$ est laissé en blanc sur la page.}
\[
\xi = \mathrm{pr}_1^{*}(\alpha\eta_p) + \mathrm{pr}_2^{*}(\eta_q)
\ \in \mathrm{Ext}^1(\mathbb{Z}/p \times \mathbb{Z}/q, \mathbb{Z})
\simeq \mathbb{Z}/p \oplus \mathbb{Z}/q .
\]
 Soit
$\mu = \mathrm{ppcm}(p, q)$ et $\delta = (\delta_p, \delta_q) : \mathbb{Z}/\mu
\hookrightarrow \mathbb{Z}/p \times \mathbb{Z}/q$ le couple des réductions. En
identifiant $\mathrm{Ext}^1(\mathbb{Z}/n, \mathbb{Z})$ au dual
$\mathrm{Hom}(\mathbb{Z}/n, \mathbb{Q}/\mathbb{Z})$, $\eta_n$ correspondant à
$1 \mapsto \frac1n$, on trouve
\[
\delta_p^{*}(\eta_p) = q'\,\eta_\mu, \qquad \delta_q^{*}(\eta_q) = p'\,\eta_\mu,
\qquad q' = \frac{\mu}{p},\ p' = \frac{\mu}{q} ,
\]
et donc la classe de l'extension $\widetilde{\Pi}'_0 = \delta^{*}\widetilde{\Pi}_0(q)$
de $\mathbb{Z}/\mu$ par $\mathbb{Z}$ :
\[
\xi' = \delta^{*}\xi = (p' + q'\alpha)\,\eta_\mu .
\]

\pagerange{34}{34}
\subsection*{Le cas $p = 6$, $q = 4$ (page 34)}

Pour $p = 6$, $q = 4$, on a $\mu = 12$, $p' = 3$, $q' = 2$, et avec $\alpha = -1$,
\[
\xi' = (3 - 2)\,\eta_{12} = \eta_{12} \,!
\]
L'extension obtenue au-dessus de $\mathbb{Z}/12$ est l'extension canonique : c'est
celle qui, tirée en arrière par l'abélianisation, donne $S\mathcal{T}_{1,1}$ à
partir de $\mathcal{T}_{1,1}$ (page 17). Le $\mathbb{Z}/6$ est celui du revêtement
défini par $\mathcal{O}(1)$, le $\mathbb{Z}/4$ celui du $\mathbb{D}'_3$.\footnote{La
page écrit « $\alpha = -1$ ? » avec un point d'interrogation ; le signe est
fixé par le caractère qui suit.}

Ce caractère est
\[
\chi : \Pi^{\mathbb{D}}_{0,3} \longrightarrow \mathbb{Z}/6, \qquad
\rho \mapsto -2, \quad \sigma \mapsto 3, \quad \varepsilon_0 = \sigma\rho \mapsto 1,
\quad \rho_0 = \varepsilon_0^{\,2} \mapsto 2 ,
\]
défini directement par l'action de $\Pi^{\mathbb{D}}_{0,3}$ sur le revêtement
principal de $(U_{0,3}, \mathfrak{S}_3)$ de groupe $\mu_6 \simeq \mathbb{Z}/6$ qui
réduit le $\mathbb{G}_m$-torseur $PU_{0,3} = \mathcal{O}_{U_{0,3}}(1)^{*}$ (page 22).
Comme $PU_{0,3}$ se déduit de ce $\mu_6$-torseur par l'inclusion
$\mu_6 \subset \mathbb{G}_m$, dont la suite d'homotopie est
$\mathbb{Z} \xrightarrow{6} \mathbb{Z} \to \mathbb{Z}/6$, on a un carré cartésien
\[
\pi_1(PU_{0,3}, \mathfrak{S}_3) = \widetilde{\Pi}{}^{\mathbb{D}}_{0,3} \simeq
\Pi^{\mathbb{D}}_{0,3} \times_{\mathbb{Z}/6} \mathbb{Z} ,
\]
le noyau $\mathbb{Z}$ étant engendré par le lacet $\varpi'$ de la fibre
$\mathbb{G}_m$. Dans ce produit fibré, $\tilde\rho = (\rho, -2)$ et
$\tilde\sigma = (\sigma, -3)$ vérifient $\tilde\rho^{\,3} = \tilde\sigma^2 = (1, -6)$,
générateur du noyau, et
\[
\pi_1(PU_{0,3}, \mathfrak{S}_3) = \langle \tilde\rho, \tilde\sigma \mid
\tilde\rho^{\,3} = \tilde\sigma^2 \rangle \simeq \widetilde{SL}(2, \mathbb{Z}) :
\]
ce sont les relèvements normalisés de la page 11 pour $n = 3$, $\nu = -1$. Ce
groupe est le groupe des tresses $B_3$, et l'espace
$(PU_{0,3}, \mathfrak{S}_3)$ le complémentaire de la cubique cuspidale
$\gamma_2^3 = 2\gamma_3^2$ de la page 22.\footnote{La page écrit
$\tilde\rho^{\,3} = \tilde\sigma^2 = \varpi'$ ; avec la convention
$(1, 6) = \varpi'$ c'est $\varpi'^{-1}$, et le signe dépend de l'orientation
choisie sur la fibre. Elle entoure l'étiquette « $-\alpha'$ » de la flèche
$\Pi^{D}_{0,3} \to \mathbb{Z}/6$ en haut de la page et « $\alpha'$ » dans le
diagramme plus bas : les deux caractères diffèrent par le signe, et chacun
engendre $\mathrm{Hom}(\Pi^{\mathbb{D}}_{0,3}, \mathbb{Z}/6)$. La
dernière phrase du corps du texte relie la page 22 à la page 34 ; ce
rapprochement est de la présente lecture.}

Enfin la page écrit, au-dessus de $\Pi^{\mathbb{D}}_{0,3} \times \mathbb{Z}/4$,
l'extension $\widetilde{\Pi}{}^{\mathbb{D}}_{0,3}(4) = \pi_1(PU_{0,3},
\mathfrak{S}_3 \times \mu_4)$ de la page 28, et, en dessous, l'image inverse
$\widetilde{\mathcal{T}}_{1,1}$ du sous-groupe $\mathcal{T}_{1,1} =
\Pi^{\mathbb{D}}_{0,3} \times_{\mathbb{Z}/2} \mathbb{Z}/4$ : une extension de
$\mathcal{T}_{1,1}$ par le même $\mathbb{Z}$, sous-groupe d'indice $2$ de
$\widetilde{\Pi}{}^{\mathbb{D}}_{0,3}(4)$.\footnote{La page relie
$\widetilde{\Pi}{}^{D}_{03}(4)$ et $\widetilde{\mathcal{T}}_{11}$ par un signe $=$
vertical ; comme $\mathcal{T}_{1,1}$ est d'indice $2$ dans $\Pi^{\mathbb{D}}_{0,3}
\times \mathbb{Z}/4$ et que les deux extensions ont le même noyau, ce ne peut être
qu'une inclusion d'indice $2$. Elle écrit aussi
$\mathcal{T}_{1,1} = \Pi^{D}_{0,3} \times_{\mathbb{Z}/12\mathbb{Z}}
\mathbb{Z}/4\mathbb{Z}$ ; voir la note sur la page 19.}

\pagerange{33}{35}
\section*{La réponse : $S\mathcal{T}_{1,1} \simeq S\Pi^{\mathbb{D}}_{0,3}$ (pages 33 et 35)}

\subsection*{Page 33}

La question de la page 19 reçoit sa réponse :
\[
S\mathcal{T}_{1,1} \simeq \widetilde{SL}(2, \mathbb{Z})
\simeq \langle \rho, \sigma \mid \rho^3 = \sigma^2 \rangle
\simeq \Pi^{\mathbb{D}}_{0,3} \times_{\mathbb{Z}/6} \mathbb{Z}
= S\Pi^{\mathbb{D}}_{0,3} ,
\]
par $\rho \mapsto (\rho, 2)$, $\sigma \mapsto (\sigma, 3)$. En effet, cette
application envoie $\omega = \rho^3 = \sigma^2$ sur $(1, 6)$, générateur du noyau
du produit fibré ; elle est surjective sur $\Pi^{\mathbb{D}}_{0,3}$, de noyau
$\langle \omega \rangle$, et induit un isomorphisme de ce noyau sur
$6\mathbb{Z}$.\footnote{La page commence par un « 4) » biffé
(« $\widetilde{SL}(2, \mathbb{Z}) \simeq \Pi \ldots$ »), puis un « 1) ». La
vérification est ajoutée. C'est aussi le cas $n = 3$ de la proposition de la page
11, avec l'isomorphisme opposé $\rho \mapsto 2$ au lieu de $\rho \mapsto 4$ sur
$\mathbb{Z}/6$.} La suite de la page, tête-bêche, reprend le diagramme de la
page 30, avec la question de savoir comment $\widetilde{\Pi}$ s'envoie dans
$\widetilde{\Pi}(4)$ au-dessus de $\Pi'$.\footnote{La flèche verticale porte un
« ? » et un mot illisible. Un mot bref d'une autre main, à l'encre rouge, daté du
27 novembre 1981, occupe aussi la page ; c'est la seule date écrite sur ces
feuillets.}

\pagerange{35}{35}
\subsection*{Page 35}

Dans $S\Pi^{\mathbb{D}}_{0,3} = \langle \rho, \sigma \mid \rho^3 = \sigma^2 \rangle$,
soit $\omega = \rho^3 = \sigma^2$ et $\omega' = \omega^{-1}$, générateur du noyau.
La somme amalgamée le long de $\omega' \mapsto \varpi^4$ est
\[
S\Pi^{\mathbb{D}}_{0,3}(4) = \langle \rho, \sigma, \varpi \mid \varpi
\text{ central},\ \rho^3 = \sigma^2 = \varpi^{-4} \rangle ,
\qquad
1 \to \mathbb{Z}\varpi^4 \to S\Pi^{\mathbb{D}}_{0,3}(4) \to
\Pi^{\mathbb{D}}_{0,3} \times \mathbb{Z}/4 \to 1 ,
\]
avec $\varpi \mapsto (1, 1)$ ; la page l'identifie à
$\pi_1(PU_{0,3}, \mathfrak{S}_3 \times \mu_4)$. L'application
\[
S\mathcal{T}_{1,1} \longrightarrow S\Pi^{\mathbb{D}}_{0,3}(4), \qquad
\rho \mapsto \rho\,\varpi^2, \quad \sigma \mapsto \sigma\,\varpi^3, \quad
\omega \mapsto \varpi^2, \quad \omega^2 \mapsto \varpi^4 ,
\]
est bien définie, relève $\mathcal{T}_{1,1} \hookrightarrow \Pi^{\mathbb{D}}_{0,3}
\times \mathbb{Z}/4$, $\rho \mapsto (\rho, 2)$, $\sigma \mapsto (\sigma, 3)$, et
envoie le noyau $\mathbb{Z}\omega^2$ isomorphiquement sur $\mathbb{Z}\varpi^4$ :

\begin{tikzcd}[column sep=small, row sep=small, nodes={font=\scriptsize}]
1 \arrow[r] & \mathbb{Z}\varpi^4 \arrow[r] & S\Pi^{\mathbb{D}}_{0,3}(4) \arrow[r] & \Pi^{\mathbb{D}}_{0,3} \times \mathbb{Z}/4 \arrow[r] & 1 \\
1 \arrow[r] & \mathbb{Z}\omega^2 \arrow[r] \arrow[u] & S\mathcal{T}_{1,1} \arrow[r] \arrow[u] & \mathcal{T}_{1,1} \arrow[r] \arrow[u, hook] & 1
\end{tikzcd}

Donc $S\mathcal{T}_{1,1}$ est exactement l'image inverse de $\mathcal{T}_{1,1}$
dans $S\Pi^{\mathbb{D}}_{0,3}(4)$, sous-groupe d'indice $2$, et c'est la réponse
au point d'interrogation de la page 28.\footnote{La page écrit la relation
« $\rho^3 = \sigma^2 = \dot\varpi^4$ » avec un « NB $\dot\varpi^4 =
\ldots^{-1} = \omega'$ », le point au-dessus de $\varpi$ étant incertain. Seule la
convention $\rho^3 = \sigma^2 = \varpi^{-4}$ donne les valeurs
$\omega \mapsto \varpi^2$, $\omega^2 \mapsto \varpi^4$ inscrites sur la flèche :
avec $\varpi^{+4}$, $\omega$ irait sur $\varpi^{10}$ et l'image serait d'indice
$5$ dans l'image inverse de $\mathcal{T}_{1,1}$. La page marque aussi
« $\omega' = \varpi^{-1}$ » au-dessus du noyau de la première ligne, où $\varpi$
désigne vraisemblablement $\omega$. Au bas de la page, tête-bêche :
« $\pi_1(U_{03}, \mathfrak{S}_3 \times \mu_4)$ », où l'on attend $PU_{03}$ comme à
la page 34. Le rapprochement avec la page 28 est de la présente lecture.} En
termes d'espaces, ce que les pages 26, 34 et 35 assemblent sans l'écrire en une
ligne est que $S\mathcal{T}_{1,1}$ est le groupe fondamental du champ
$(PU_{0,3}, \mathbb{D}'_3)$, où $\mathbb{D}'_3 = \mathfrak{S}_3
\times_{\mathbb{Z}/2} \mu_4 \subset \mathfrak{S}_3 \times \mu_4$ agit sur le
$\mathbb{G}_m$-torseur $\mathcal{O}(1)^{*}$ — c'est-à-dire, si l'on suit la page
26, du torseur $PM_{1,1}$ au-dessus de $M_{1,1}$.\footnote{Cette formulation est
de la présente lecture ; elle découle des énoncés des pages 34 et 35, mais aucune
page ne l'écrit.}

\end{document}
